% !TeX program = pdflatex
% ALIUS Bulletin standalone facsimile source.
% Compile from the ALIUS-bulletin project root for exact visual fidelity.
% If the shared PDF asset is absent, this file falls back to extracted text.
\providecommand{\ALIUSRootPrefix}{}

\ifdefined\ALIUSIssueBuild
  \IfFileExists{\ALIUSRootPrefix Shared-assets/original-pdfs/Issue03/alius-bulletin-issue-03-winkelman-evolutionary-neuroanthropology-of-consciousness-exploring-the-diversity.pdf}{%
    \includepdf[pages=-,fitpaper=true]{\ALIUSRootPrefix Shared-assets/original-pdfs/Issue03/alius-bulletin-issue-03-winkelman-evolutionary-neuroanthropology-of-consciousness-exploring-the-diversity.pdf}%
  }{%
    \typeout{ALIUS warning: missing source PDF Shared-assets/original-pdfs/Issue03/alius-bulletin-issue-03-winkelman-evolutionary-neuroanthropology-of-consciousness-exploring-the-diversity.pdf}%
  }%
  \expandafter\endinput
\fi

\documentclass[a4paper,11pt]{article}
\usepackage[margin=22mm]{geometry}
\usepackage{pdfpages}
\usepackage[hidelinks]{hyperref}
\usepackage{microtype}
\usepackage{parskip}
\title{The evolutionary neuroanthropology of consciousness: Exploring the diversity of conscious states across cultures}
\author{Michael Winkelman Interviewed by Martin E. Fortier}
\date{ALIUS Bulletin 3}
\begin{document}
\IfFileExists{\ALIUSRootPrefix Shared-assets/original-pdfs/Issue03/alius-bulletin-issue-03-winkelman-evolutionary-neuroanthropology-of-consciousness-exploring-the-diversity.pdf}{%
  \includepdf[pages=-,fitpaper=true]{\ALIUSRootPrefix Shared-assets/original-pdfs/Issue03/alius-bulletin-issue-03-winkelman-evolutionary-neuroanthropology-of-consciousness-exploring-the-diversity.pdf}%
}{%
  \maketitle
  \section*{Metadata}
  \textbf{Piece type:} interview\par
  \textbf{DOI:} 10.34700/krg3-zk35\par
  \textbf{Keywords:} altered states of consciousness, enculturation, evolutionary neuroanthropology, hallucinogens, shamanism\par
  \section*{Abstract}
  In this interview, Michael Winkelman and Martin Fortier discuss the extent to which consciousness is grounded in deep evolutionary mechanisms and can be enculturated. First, the main tenets of two neuroanthropological approaches to consciousness and culture are outlined. Next, the upsides and downsides of evolutionary psychology are examined; the fruitfulness of this approach in the study of cultural phenomena such as shamanism is debated. The authors then discuss the promises of the "big data" approach to the study of religion as well as evolutionary puzzles about religion. Turning to issues bearing on the taxonomy of consciousness, the interview explores how consciousness should be individuated and especially how many "modes" of consciousness should be identified based on what we know of the biology and phenomenology of altered consciousness. Winkelman's concepts of the "integrative mode of consciousness" and "psychointegrators" are subsequently examined. Next, the interview addresses both how alterations of consciousness are universally similar (the perennialist view) and can also be enculturated (the constructivist view). Finally, the authors discuss issues around the cultural use of hallucinogens (a.k.a. visionary plants): what is the best method to study them? And for how long have humans used them?\par
  \section*{Extracted Text}
  \begingroup
  \small
  \noindent ALIUS Bulletin n3 (2019) aliusresearch.org/bulletin The evolutionary neuroanthropology of consciousness Exploring the diversity of conscious states across cultures An interview with Michael Winkelman by Martin E. Fortier Michael Winkelman michaeljwinkelman@gmail.com (Retired) School of Human Evolution and Social Change Arizona State University Tempe, AZ, USA Martin E. Fortier martin.fortier@ens.fr Institut Jean Nicod ENS/EHESS, Paris, France Citation: Winkelman, M. \& Fortier, M.E.. (2019). The evolutionary neuroanthropology of consciousness : exploring the diversity of conscious states across cultures. An interview with Michael Winkelman. ALIUS Bulletin, 3, 45-97. doi: 10.34700/krg3-zk35 Abstract In this interview, Michael Winkelman and Martin Fortier discuss the extent to which consciousness is grounded in deep evolutionary mechanisms and can be enculturated. First, the main tenets of two neuroanthropological approaches to consciousness and culture are outlined. Next, the upsides and downsides of evolutionary psychology are examined; the fruitfulness of this approach in the study of cultural phenomena such as shamanism is debated. The authors then discuss the promises of the "big data" approach to the study of religion as well as evolutionary puzzles about religion. Turning to issues bearing on the taxonomy of consciousness, the interview explores how consciousness should be individuated and especially how many "modes" of consciousness should be identified based on what we know of the biology and phenomenology of altered consciousness. Winkelman's concepts of the "integrative mode of consciousness" and "psychointegrators" are subsequently examined. Next, the interview addresses both how alterations of consciousness are universally similar (the perennialist view) and can also be enculturated (the constructivist view). Finally, the authors discuss issues around the cultural use of hallucinogens (a.k.a. visionary plants): what is the best method to study them? And for how long have humans used them? keywords: altered states of consciousness, enculturation, evolutionary neuroanthropology, hallucinogens, shamanism. Your work is very multidisciplinary and spans several fields (at least the following: cultural anthropology, neuroanthropology, neurotheology, cognitive science of religion, evolutionary anthropology, medical anthropology and evolutionary psychology). What was your training like? Who were the most important inspirational figures in your intellectual development? I wandered from psychology to anthropology to cross-cultural psychology and cross- cultural studies (ethnology, holocultural studies) to Buddhist psychology, then the neurosciences and neurophenomenology and evolutionary psychology in the search Michael Winkelman - The evolutionary neuroanthropology of consciousness 46 ALIUS Bulletin n3 (2019) aliusresearch.org/bulletin of a method for understanding the biological bases of spiritual phenomena. My training was eclectic, largely self-driven; I went to graduate school at the School of Social Sciences, University of California, but most of my coursework was independent study. I was in search of the next paradigm, so the current paradigms seemed little relevant. But there were some important figures along the way. Intellectually, I was inspired by the work of Charles Laughlin, especially his co- authored Brain, Symbol and Experience (Laughlin, McManus, \& D'Aquili, 1990) which set the foundation for the neurophenomenological approaches. Unfortunately, I did not grasp the full significance of Laughlin's co-authored The Spectrum of Ritual (D'Aquili, Laughlin, \& McManus, 1979) when I first read it, but this too eventually influenced my work. Most of your research work tackles topics in neuroanthropology. Importantly, the field of neuroanthropology is split, as it were, between two distinct schools. The first coincides with the biogenetic structuralist approach as defined by the seminal work of Charles Laughlin and colleagues (Laughlin \& D'Aquili, 1974; Laughlin et al., 1990). According to this school, the brain is endowed with many innate and universal neural (or "neurognostic") structures that explain behaviors and beliefs observed across the world. The second school is more recent and is best epitomized by the volume edited by Daniel Lende \& Greg Downey (2012), The encultured brain, and by work in cultural neuroscience conducted by Joan Chiao and colleagues (2016). The work of Andreas Roepstorff and coworkers (2010) also belongs to this approach. In contrast to the first school of neuroanthropology, the second one stresses the plasticity of the brain and argues that many neural structures can be shaped and influenced by idiosyncratic cultural patterns. It seems that your own work belongs to the first school of neuroanthropology. What makes you think the human brain is best described as endowed with innate and fixed structures rather than plastic ones? Your writings speak very highly of the biogenetic structuralist school, but what is your view on the "culturalist" school of neuroanthropology? Is your theoretical approach completely at odds with the latter or do you share at least some of their views? First I would say that the cross-cultural principles, even universals, of magico- religious practices, especially shamanism, speak to some underlying biological factors that produce these similarities. Certainly the physical and social environment provide influences, but I think the notion of cross-culturally distributed principles of shamanism, religion, meditation and spirituality speaks strongly to the underlying biological bases as the structural foundations. In terms of the neuroanthropology traditions, my work has been largely focused within the first school, understanding religious and spiritual universals within the biogenetic structuralist and neurognostic approaches. But I also think it abundantly clear that the brain is a highly plastic organ. At the same time, what is subject to Michael Winkelman - The evolutionary neuroanthropology of consciousness 47 ALIUS Bulletin n3 (2019) aliusresearch.org/bulletin plasticity are the innate tendencies, the extent to which they are elicited, suppressed or developed. So to me they are not opposing ideas, but different points of departure. The innate modular structures are part of our deep primate heritage, their slow manifestations across the primate and hominid and hominin lines speaks to accretion. While I have primarily focused on the first neuroanthropology approach as a tool to explain universals, my own work has also addressed this notion of cultural and social effects on plasticity. This is exemplified in my research (for a summary, see: Winkelman, 2018b) on the sociophysiological and psychological dynamics of possession states. In this context I have made my clearest statements regarding the roles of social circumstances in shaping plasticity of the brain into a very specific alteration of consciousness found cross-culturally and subsumed under various forms of possession. Some of my earliest research reflects concerns with the social effects on plasticity of innate responses. In a study of the effects of formal education on extra sensory perception and of the effects of socialization on the manifestations of psi (Winkelman, 1982), I address issues related to how socialization affects experience and awareness of the world, very much concerned with the plasticity of innate capacities. So this plasticity is why cultural spiritual traditions are so important. They can take various disturbances in life, such as illness, injury, abandonment, being orphaned, etc. and shape the body-mind response, using the consequences of these disturbances to alter the overall normal cognitive-emotional developmental trajectory of the person. Such disturbances, especially illness, are used in the traditions of shamanism and meditation, taking the disturbance in normal development to enable one to engage our innate structures and develop and associate them in different ways than that normally intended by the unfolding of our nature and the cultural evolution of our varied capacities. The shaman is a master in engaging, combining, integrating the various modular structures, while the meditator is adept at stimulating, totally isolating and disenabling innate mental structures, leading to experiences of pure light, love and joy, or void/nothingness. A good example of your biogenetic structuralist take on various cultural phenomena is your theory of shamanism. According to you, the main features of shamanism are to be understood in terms of innate neurobiological structures ("neurognostic structures") present in every human: "The cross-cultural manifestations of basic experiences related to shamanism (e.g., soul flight, death-and-rebirth, animal identities) illustrates that these practices are not strictly cultural but are structured by underlying, biologically inherent structures. These are neurobiological structures of knowing that provide the universal aspects of the human brain/mind" (Winkelman, 2010a, p. 38; and also: Winkelman, 2002b, 2002a). Yet, drawing upon Ake Hultkrantz's work you also acknowledge that shamanism arises out of environmental adaptations: "The worldwide Michael Winkelman - The evolutionary neuroanthropology of consciousness 48 ALIUS Bulletin n3 (2019) aliusresearch.org/bulletin similarity in shamans derives from the psychobiological bases of human consciousness and its adaptation to social and ecological conditions of hunter-gatherer societies" (Winkelman, 2010a, p. 64). Regarding the study of shamanism, is the model you endorse one in which neurognostic structures directly cause specific cultural traits to emerge, or one with an interactionist twist where cultural traits emerge from the interplay of neurognostic and environmental constraints? All of genetics and innate capacities unfold in interaction with the environment. I see the foraging band structure with a fission-fusion dynamic as the environment in which shamanism emerged, something deep in hominin prehistory (Winkelman, 2010b). These ideas about the commonalities in chimp and hominin rituals emerged in my work circa with the second edition of Shamanism (2010a). These ancient ritual structures reflect what is at issue in the first statements you review. The second issue about the effects of foraging environments on the manifestation of shamanism is more related to the contrast with other kinds of magico-religious practitioners (Winkelman, 1990). I think that many of these same capacities at the core of shamanism are manifested in the experiences and behavior of possession cult mediums, whose distinctive features emerged from the formative influences exerted by the oppressive societal dynamics of social stratification and political integration, probably exacerbated by patrilineal social structures and compromised nutrition. So I think that the traits emerge both from the innate tendencies, as well as environmental provocation. For instance, the altered state of consciousness (ASC) associated with extreme fasting also emerge "naturally" from the effects of food shortage and short-term starvation. I think the evidence indicates that out-of-body experiences can be provoked by many stressors, including accidents, long-distance running, etc. I emphasize these in the sense of a neurognostic structure, something innate about how we perceive the world. The cultural part is what we believe about these experiences. Such cultural beliefs may have many sources apart from the neurognostic dynamics produced by the innate dynamics of these experiences. " I think that the traits emerge both from the innate tendencies, as well as environmental provocation. This distinction of neurognostic versus cultural structures raises the question of the extent to which shamanic traditions depend on culture. Clearly cultural loss leads to loss of shamanic practices, but whether this is simply the consequence of some loss of "knowledge" or some broader consequences of cultural disintegration is not " Michael Winkelman - The evolutionary neuroanthropology of consciousness 49 ALIUS Bulletin n3 (2019) aliusresearch.org/bulletin clear. To me, the remarkable similarities of shamanic practices in societies around the world speaks to something innate about the fundamental dynamic of shamanism, something that can emerge from innate cognitive, social and emotional tendencies rather than something sustained primarily through cultural belief; this is supported by converging evidence from ethology, neuropsychology and cross- cultural studies (Winkelman, 2010c). However, when these tendencies are provoked in the modern world, their sequela and consequences are quite different from that seen under conditions of foraging societies because the formative social influences (and cultural beliefs) are dramatically different. The environment is also a factor here, where animals are clearly central to adaptation and survival, whether as food or predators, or even as sentinels. I think that a core part of shamanism was a relationship with the world of animals on so many levels. This gets reduced to cows and chickens and later dogs and cats. Not as complex animal-based environmental relations that creates identities and powers and channels into nature typical of foraging societies. According to you, shamanic concepts such as the belief in a "traveling soul" or in "soul flight" are to be explained by innate modules (Winkelman, 2002b, 2002a, 2010a, 2017b). The concept of a module has been largely debated and what is meant by it in Fodor's seminal piece (Fodor, 1983) is very different from what is meant in contemporary evolutionary psychology (C. Barrett \& Kurzban, 2006). What definition of a module do you endorse? For example, do you take automaticity and informational encapsulation to be definitional of what a module is? Moreover, do you think with Fodor that the modularity of the mind is restricted to perception and does not concern cognition (reasoning, concepts, etc.); or, along with Tooby \& Cosmides (1992), Sperber (1996a) and others, do you champion the massive modularity hypothesis (the view that the mind is modular throughout)? My conceptualization of innate modules is primarily informed by Gardner's work, as well as the broader concept of innate operators used by d'Aquili to conceptualize universals of spiritual experience. These innate modules or operators are far more than perceptual devices, but engage full-blown cognitive processes providing the basis for innate intelligences. I see shamanism as exploiting a variety of these innate intelligences, and furthermore using ritual processes to combine innate intelligences, for instance in animal allies that combine the interpersonal and naturalistic intelligences. This integration of (1) interpersonal intelligence-a capacity for a "theory of mind" to infer others' mental processes and to work effectively with others through an understanding of their motivations and intentions; with (2) the naturalist intelligence that provides "expertise in the recognition and classification of the numerous species-the flora and fauna-of his or her environment" (Gardner, 2000, p. 36). This cross-modular integration is in contrast to the normal isolated function of modules. It also provides a new capacity Michael Winkelman - The evolutionary neuroanthropology of consciousness 50 ALIUS Bulletin n3 (2019) aliusresearch.org/bulletin by taking the ability to assign identity to/recognize a species and to identify relations between and among species and create a new form and level of symbolic thought for personal and social differentiation (i.e., totemism). I would endorse a massive modularity concept, one that would include the notion of several innate forms of self, following the concepts of Damasio; I think that these kinds of innate conceptual structures not only underlie the body self-experienced in out-of-body experiences, but various concepts of embodied spirits, ancestor spirits, and external spirit entities (see Winkelman, 2018a you discuss below). I think these innate intelligences may constitute forms of narrative production, such as Gardner's concepts of: -A spiritual intelligence, manifested in "a desire to know about experiences and cosmic entities that are not readily apprehended in a material sense" (Gardner, 2000, p. 40) and engaging with spiritual, noetic and transcendent experiences, inducing charisma, and using it to instill a quest for this spiritual awareness; and -An existential intelligence that reflects the cognitive aspects manifested in the spiritual intelligence, "an ability to locate oneself with respect to the furthest reaches of the cosmos [...] the significance of life, the meaning of death [...] a concern with cosmic issues" (Gardner, 2000, p. 44). When you discuss the belief in soul flight the output of the modules you are referring to is experiential-rather than conceptual or doxastic: it consists of out-of-body experiences, near death experiences, astral projection experiences, etc. (e.g., Winkelman, 2010a, p. 117). Now, you argue that these experiences somewhat automatically trigger beliefs whose content largely depends on the content of experience. This line of reasoning looks very much like that of Tylor (1871) when he suggests that beliefs in souls are triggered by altered states of consciousness such as dreaming. Yet, the relationship between experience and beliefs (or concepts) is far from simple and obvious. Unless one endorses the view that perceiving is believing-this view seems wrong to me, since one can perceive things without taking them to be real (Dokic \& Martin, 2012; Fortier, 2018a, 2018b)-it is obvious that some dissociation is to be found between experience and belief: (1) some people will believe in soul flight without experiencing any out-of- body experience; (2) and some will experience out-of-body experiences without believing in soul flights. Let me illustrate both (1) and (2). (1) During my fieldwork in the Middle Ucayali (Peruvian Amazon), a Shipibo friend of mine once told me about the experience he had of losing his soul (kaya). While my friend was quietly paddling, his canoe suddenly capsized and he fell into the river. He managed to put the canoe back in its straight position and came back home but noticed he was feeling strange: he was feeling very depressed. Intriguingly, he interpreted this feeling as his soul being gone. As a side note, it should be stressed that according to the Shipibo- Konibo model of the person, when your soul is gone, you do not become a zombie or an Michael Winkelman - The evolutionary neuroanthropology of consciousness 51 ALIUS Bulletin n3 (2019) aliusresearch.org/bulletin unconscious corpse; you still have some components active in your body-like shinan (thought and emotion) or bero yoshin (the spirit of the eye)-that enable you to feel and think. My friend then talked to a shaman and they decided to organize an ayahuasca ceremony. During that ceremony, the shaman asked my friend's soul to come back. The day after, when my friend woke up, he was not feeling depressed anymore. He took it to mean that his soul had eventually come back in his body (yora). This anecdote is very interesting because it nicely illustrates how the belief in soul travel may have nothing to do with out-of-body experiences-or indeed any anomalous experience. It can merely be a matter of feeling depressed or a bit strange. (2) If at least some people in the world seem to believe in soul flight without associating it with altered states of consciousness, conversely, some people experience altered states of consciousness without believing in soul flight. A good example of this is provided by volunteers who participate in experiments where neuroscientists purposely induce out-of-body experiences (e.g., Blanke, Landis, Spinelli, \& Seeck, 2004). Importantly, in spite of having those out-of-body experiences, participants are mostly secular Westerners who do not believe in soul flight neither before nor after the experience. The out-of-body "module" is indeed triggered but no belief in soul flight is formed as a result of it. As these two examples demonstrate, it is not obvious that religious beliefs or concepts can be explained strictly by the activation of modules whose output is experiential: such beliefs can be held without the module being activated; and the module can be triggered without such beliefs being formed. What is your take on this? How do you conceive of the relationship between experiential modules and beliefs? I think that it is important to make the distinction between, on one hand, the structural and functional aspects of the experience, and on the other hand the explanatory, epistemic and metaphysical constructs to provide an account of the experience. The former is where we see the innate structures manifested; the latter metaphysical concepts may be very neurognostic (Winkelman, 2013) or they may be formed by a variety of cultural traditions and assumptions. This would include the assumptions of modern medicine and reductionist biology that view certain aspects of experience as illusory because they hold deep metaphysical conflicts with the assumptions of biomedicine (for the cultural aspects of biomedicine see: Winkelman, 2009, Chapter 5). I think that Tylor was addressing the obvious in attributing the origins of supernatural beliefs to experiences of dreams, especially various forms of lucid and precognitive dreams. Humans have experiences manifested cross-culturally because they have something to do with how our biology operates-or malfunctions. There are certainly many experiences humans have that imply a dualistic reality, especially those involving out-of-body experiences (OBEs) that can be occasioned by diverse means. The fact that OBEs can be caused by such diverse conditions-various medical ailments, psychological manipulation, various drugs, drumming, long- Michael Winkelman - The evolutionary neuroanthropology of consciousness 52 ALIUS Bulletin n3 (2019) aliusresearch.org/bulletin distance running, etc.-indicates we are seeing something about the structural features of the organism's perceptual/conceptual capacities. I see in the OBE first the disassembling of the normal integration of different modular structures-visual field, body sense, egoic identity-and then subsequently the emergence of a deep and ancient self-structure, perhaps related to the mimetic mind (Winkelman, 2017b) or one of the non-verbal forms of self identified by Damasio. How we understand-even perceive and relate to-these experiences depends on our prior orientation. In the context of an animistic pre-modern worldview the spiritual dimensions and explanations of these experiences is natural and predominant. When you take post-modern science students into a lab, of course their explanations will be different. But furthermore don't think for a moment that those who "participate in experiments where neuroscientists purposely induce out-of-body experiences" are really experiencing the same thing as a shaman who just collapsed after drumming and dancing for six hours, or someone lost in an ayahuasca journey, or someone clinically dead and laying on the side of the rode with paramedics working to stimulate a heartbeat. There are many dimensions to the dynamics of the OBE besides some simple phenomenological descriptions about body and self-awareness and identity. It is not entirely clear to me what the explanatory gain is to resort to psycho- evolutionary accounts of cognitive/cultural phenomena. To illustrate my skepticism, let me take the example of a psycho-evolutionary account of the fusiform face area (FFA). The FFA has been discussed in evolutionary terms both by yourself (Winkelman, 2018a, p. 6) and by other evolutionary psychologists/cognitive anthropologists (Sperber \& Hirschfeld, 2004, pp. 40-42). Specifically, in "An ontology of psychedelic entity experiences in evolutionary psychology and neurophenomenology," you propose that the reason why people hallucinate many eyes and faces under the effect of psychedelic compounds is that their FFA becomes hyper-activated (Winkelman, 2018a, p. 7). But what is the actual "added value" of arguing that the FFA has been evolutionarily-shaped? Some evidence suggests the FFA is in fact not evolutionarily shaped, but depends on expertise: it is an expertise-dedicated area (Gauthier \& Nelson, 2001). This explains why experts in nonface objects (e.g., cars) have the FFA activated by objects other than faces. If Gauthier's theory of the FFA is correct, this means that those who have the FFA being activated by faces are "face experts" and that their FFA "module" has been shaped by expertise rather evolution. The point worth stressing is that one can still say (a) "hallucinations featuring multiple faces correlate with the hyper-activity of the FFA," even if the FFA turns out to be shaped through development-rather than through evolution. This is what I mean when I say that evolutionary explanations bring no "explanatory added value." Another way to put this objection would be to take the example of cars. Gauthier has shown that the FFA could become a "module" specialized in car detection. Now, no one would claim that we have a car module that has been shaped through evolution because obviously cars did not exist in the Pleistocene. Yet, it is still possible to propose that (b) "hallucinations Michael Winkelman - The evolutionary neuroanthropology of consciousness 53 ALIUS Bulletin n3 (2019) aliusresearch.org/bulletin featuring multiple cars correlate with the hyper-activity of the FFA (or any other brain area(s) specialized in the detection of cars)." At the end of the day, it seems that the explanatory merit of statements such as (a) or (b) lie not in the appeal to some putative evolutionary origin-neither (a) nor (b) do such a thing-but simply in reducing phenomenology to some brain function. In other words, the explanatory added-value of such statements comes from reducing cognitive functions to brain activity. What is true of the FFA is true of most-if not of all-modules. Their evolutionary origin is very controversial and some authors have shown that modules can be accounted for by developmental mechanisms rather than evolution (Elman et al., 1996; Karmiloff- Smith, 1992, 2009). When you argue that the "social cognition module," the "biological module," the "physical module," etc. explain crucial aspects of shamanism, the claim seems to be that your evolutionary account sheds some light on the phenomenon at hand. But it can be objected that the same could be said even if all our modules have been shaped by development and have nothing to do with evolution. For example, you can say "people anthropomorphize their surroundings because of the social cognition module", even if it turns out this module is not an adapted function (i.e., is not shaped by evolution); even if evolutionary processes are not involved, it will always remain that we have cognitive functions specialized in processing social stimuli and that anthropomorphism can always be attributed to them. In short, according to you, what is the real explanatory gain of evolutionary accounts of cultural phenomena? I think that the evolution of function is what is the focus of selection, and that this always happens through repurposing old hardware. So the FFA undoubtedly faced many selective pressures, one of which is its detection of faces. The point here about the innate modules is that they were designed to provide a certain functional response. My approach is not to say that the responses are just about the brain- obviously they are stimulated by social and environmental context. New function can emerge out of new connections rather than new hardware. The importance of the innate modules and evolved psychology approaches is the following. When we find recurrent patterns of experience or behavior across cultures, the traditional cultural explanation of religious or spiritual beliefs offers no explanatory power. Why is an out-of-body experience reported around the world? Not because cultures believe this, but because of innate dispositions. When we encounter universals or highly repeated phenomena across cultures, there must be a biological reason. Now notably we do not have people from cultures around the world reporting the presence of cars, not even cars staring at them (which one might even expect since the two headlights are like eyes and the grill like a big mouthed grin!). So is it as you say "[T]he explanatory added-value of such statements comes from reducing cognitive functions to brain activity"? To me, that is a basic function of neurognostic approaches in explaining religious universals. It has to do with how our brain operates. Of course there is the added question of how it is that religious, Michael Winkelman - The evolutionary neuroanthropology of consciousness 54 ALIUS Bulletin n3 (2019) aliusresearch.org/bulletin ritual and spiritual activities have the strong tendency to stimulate these innate cognitive tendencies. More to come but see: Winkelman, 2017b, 2018a. " The importance of the innate modules and evolved psychology approaches is the following. When we find recurrent patterns of experience or behavior across cultures, the traditional cultural explanation of religious or spiritual beliefs offers no explanatory power. If cultural phenomena studied by anthropologists-e.g., animistic and totemistic beliefs/practices-are in fact underpinned by universal adapted functions, one may then wonder how psycho-evolutionary explanations can make sense of the heterogeneous distribution of these phenomena. Indeed, humans are not animists or totemists everywhere in the world and to the same degree (e.g., Descola, 2013; Ingold, 2000; Testart, 2012). If cognitive adapted functions-such as the social cognition module-are universal, and if these cognitive adaptations underlie cultural phenomena such as animism and totemism (Winkelman, 2002a), then why is it that animism and totemism are not found everywhere across the world (Descola, 2013, Epilogue)? The degree and nature of manifestation of innate capacities are subject to socialization. This is very clear in the area of culture bound syndromes or culture reactive syndromes, where, for instance, innate tendencies such as the startle response may be extremely activated in some cultures to the point it creates a culturally induced illness. The manifestation of innate capacities is quite variable, but inevitable in some form or degree. Animism is variable, but as Guthrie (1993) has shown very convincingly, people everywhere are animist. Depending on your concept of totemism-but let's take "thinking in terms of animal metaphors"-people everywhere are also totemist. Animism is not expressed only in religious ideation, but in concepts of everyday life, expressed in why our cars have problems, why the computer doesn't work, why the machine works sometimes and sometimes not. Innate modules do not mean that the capacity is expressed in everyone and everywhere to the same degree. Obviously social circumstances can elicit or repress certain tendencies. But is there a culture with no animist beliefs? Are there cultures with no animal metaphors to express emotions, tendencies, social relations? I don't think so. Show me. J However, our metaphysical world is less animist and totemist, etc. than many pre- modern cultures. Why so? Many answers, but one basic one is the social environment. When you live in a culture where the traditions of your ancestors still " Michael Winkelman - The evolutionary neuroanthropology of consciousness 55 ALIUS Bulletin n3 (2019) aliusresearch.org/bulletin organize your day to day, monthly and seasonal activities, such as in agriculture, the compulsion to follow your ancestors' advice is very strong and probably adaptive. Ancestor worship is an adaptive way of maintaining vital traditions and mediumistic communication with them to get advice on how to address problems is probably an adaptation for accessing the knowledge of our collective memories. But when we live in cultures where the traditions of our ancestors, including their work skills, habits and beliefs, are totally irrelevant to our day to day adaptations, small surprise that ancestor worship is not important. Ditto for animism. I think that it is not possible to fully experience nature and perceive its animistic qualities unless one is immersed in it, solely and far from civilization. When we live in intimate relations with nature, our sense of animism also must shift. After all there are all kinds of non-human living things and natural forces to be experienced in nature. So the ways in which the innate modules are elicited, their cross-entrainments with other modules and cognitive process, and their roles in culture produce a lot of variation. Shamanism provided traditions for engaging and elevating certain aspects of consciousness adaptive for our foraging ancestors and their relations to nature. Back to the plasticity issue. " [T]he ways in which the innate modules are elicited, their cross-entrainments with other modules and cognitive process, and their roles in culture produce a lot of variation. Shamanism provided traditions for engaging and elevating certain aspects of consciousness adaptive for our foraging ancestors and their relations to nature. You have been a pioneer in using databases to address big anthropological questions. In particular, you have investigated the relationship between economic systems, social structures, and types of magico-religious practitioners (priests, shamans, healers, etc.). Rather than theorizing about these variables on the basis of a single fieldwork, as many anthropologists do, you looked at the Human Relations Area Files (HRAF), gathered the relevant data, and resorted to statistical tools to answer the questions you were interested in. This led to some groundbreaking findings (Winkelman, 1986, 1992). In recent years, several large projects have similarly tried to address big questions about religion by resorting to databases. The two most significant of them are certainly the SESHAT database (seshatdatabank.info) and the Database of Religious History (religiondatabase.org). What do you think of those recent projects? " Michael Winkelman - The evolutionary neuroanthropology of consciousness 56 ALIUS Bulletin n3 (2019) aliusresearch.org/bulletin Strikingly, the questions tackled by recent dataset-based research programs completely overlook the topic of altered states of consciousness. Do you think big data could significantly help shed new light on altered states of consciousness or are there intrinsic limitations to this method? In your recent work, you do not use databases anymore; but are you planning to use such tools again in the future? While I used some data from the HRAF for my research, my research was based on a subsample from the Standard Cross-cultural Sample. There are some subtle differences between them. I feel that it is unfortunate that anthropologists and other social scientists don't rely more on cross-cultural data to address questions regarding social and cultural universals and social evolution of phenomena, religious or otherwise. Such data is essential for more valid generalizations. So it is good to see new data sets being developed to bring empirical data to address questions about the cultural evolution of religion. Eventually we will need both synchronic and diachronic data sets, and the significance of having a pinpointed culture and time for valid synchronic analysis may not be incorporated into some efforts. How we can incorporate such different info into a single data set is challenging. From some of the earliest cross-cultural research on magico-religious practitioners (e.g., Bourguignon \& Evascu, 1977), the cross-cultural variation in the alteration of consciousness has been a significant area of research and theory. Briefly, Bourguignon was the first to report empirical data on the association of possession ASC with more complex societies. My subsequent research, building on her concepts, has shown that possession appears as the predominant socially recognized ASC under conditions of political complexity and the associated conditions of oppression of females. To me, this is one of the important contributions of cross- cultural research in general, and specifically with respect to the issue of innate modules. One of the arguments against a cross-cultural shamanism is that ASC are not expressed the same everywhere. True. But what explains the patterns of variation? This is where cross-cultural research can help us understand how a set of innate capacities, in this case for alterations of consciousness, may be expressed in a variety of forms, depending on local circumstance. My cross-cultural data base on magico-religious practitioners and altered states of consciousness is a bit of a paradox. On one hand, its initial formulation provided the basis for my dissertation, Shamans, Priests and Witches, a variety of articles, and most importantly my career of developing concepts regarding the universal biological bases of shamanism. After I finished my dissertation, I received a National Science Foundation Dissertation Improvement Grant and I completely reformulated the variables based on the experience of the dissertation, recoded the data, and did coding reliability checks. This new database was never analyzed (but is available from HRAF and on researchgate.net). For a variety of reasons my career Michael Winkelman - The evolutionary neuroanthropology of consciousness 57 ALIUS Bulletin n3 (2019) aliusresearch.org/bulletin took some different directions (neurotheology, cross-cultural relations, medical anthropology, psychedelics) and I never returned to analyze the new data, or for that matter, many of the obvious questions suggested by the original data set. Will I return to analyze the data? I think about it from time to time, but I generally feel like I am too old to want to go learn how to use another statistical package. I hope someone will discover this unanalyzed data and do something with it. Principal areas that might be addressed? (1) the evolution of priesthoods. What are the origins in ancestor worship and what are the factors that contribute to their emergence as dominant ritual structures of society epitomizing the features of religions and priests? (2) ASC and healing: Are different types of ASCs associated with different illness ideologies?; How does this data set relate to Murdock's data on Theories of Illness (Murdock, Wilson, \& Frederick, 1980)? The Cognitive Science of Religion (J. L. Barrett, 2000, 2007; Pyysiainen, 2013) is one of the most ambitious programs of naturalization of religion. Although this research program has made some important findings, it has consistently overlooked the role experience-and especially altered consciousness-plays in religion. For example, it has been argued that conscious experience is uninteresting because it has no effect on religious concepts. That is, the same religious concepts will be formed whether or not altered states of consciousness are experienced (e.g., Boyer, 2001, Chapter 9; Cohen, 2007). In your HRAF-based work on magico-religious practitioners, you have demonstrated that the type of practitioner found in a culture largely depends on modes of consciousness entertained in that culture (Winkelman, 1986, 1992). Namely, the "integrative mode of consciousness" is only found in cultures where magico-religious practitioners are shamans, healers, and mediums but not in cultures where these practitioners are sorcerers, witches, or priests. This finding seems to provide a strong argument in favor of the view that religious concepts are to a certain extent shaped by conscious experience. However, this argument works only if we can establish what the direction of causality is between magico-religious practitioner types and the integrative mode of consciousness. A first possibility is that most of the traits defining shamans, healers, and mediums are determined by the integrative mode of consciousness. But a second possibility is that most of their traits coalesce as a result of other factors (e.g., biosocial or socioeconomic factors) and subsequently cause the emergence of rituals inducing the integrative mode of consciousness. According to this account, the integrative mode of consciousness would be the effect and not the cause of the traits defining shamans, healers, and mediums and their idiosyncratic beliefs. Which of these accounts do you think is the most accurate? More broadly, your work on magico-religious practitioners posits three main variables: the biosocial function, the mode of consciousness (which is sometimes described as a biosocial function) and the Michael Winkelman - The evolutionary neuroanthropology of consciousness 58 ALIUS Bulletin n3 (2019) aliusresearch.org/bulletin socioeconomic condition. But what is the causal role of each of these variables in bringing about specific magico-religious practitioner types? Three bases of magico-religious practice reflect the fundamental impulses underlying religious life. The cross-cultural distribution of these institutionalized practices reveals that the fundamental forms of religious life are not arbitrary or simply cultural, but derived from biogenetic human impulses that are manifested across cultures and time, albeit modified in their manifestations by the reigning subsistence, social and political conditions. These cross-cultural commonalities are also manifested in the configurations of magico-religious practitioners. Societies in the sample had specific typical patterns of co-occurrence of practitioner types. 1. Practitioner: Shaman or other Healer Complex (Shaman/Healer or Healer) 2. Practitioners: Priest and Healer Complex or Medium (1 society with two Healer Complex) 3. Practitioners: Priest, Healer Complex, and Medium or Sorcerer/Witch 4. Practitioners: Priest, Healer, Medium and Sorcerer/Witch In order to reveal these underlying biogenetic functions, entailment analyses were used to identify the relationship between the formal functions of these magico- religious practitioner types and the processes for practitioner role selection (Winkelman, 1986). This revealed three major relationships between selection processes and magico-religious activities, notably the former entailing the latter rather than vice-versa. (1) Alteration of Consciousness and Healing. If there is selection for the role by (a) signs from the spirits, involuntary illness or spontaneous visions or deliberate vision quests, then (b) there is further training involving alterations of consciousness, and (c) professional activities of healing and divination. These features were typically characteristic of the practitioners of the Healer Complex (all Shamans and Shaman/Healers and some of the Healers) and the Mediums. This is a human universal, manifested and utilized differently in different types of societies and traditions. So your statement above that the integrative mode of consciousness is not found in cultures where there are Sorcerer/Witches or Priests is mistaken. All societies, including those with Priests and Sorcerer/Witches, have the integrative mode of consciousness manifested in Shamans, Shaman/Healers, Healers or Mediums. All societies have one of these shamanistic healers that exploit the potentials of the integrative mode of consciousness. Michael Winkelman - The evolutionary neuroanthropology of consciousness 59 ALIUS Bulletin n3 (2019) aliusresearch.org/bulletin (2) Political Succession and Agriculture Rites and Propitiation. If there was selection for the role through (a) some form of social succession (typically father to son) or some form of political action (i.e., political negotiations or war), then (b) the practitioner exercised political, legislative and judicial power, and (c) engaged in seasonal rituals of agricultural fertility, as well as ritual activities for propitiation of collective spirits and protection. These features were characteristic of the practitioners labeled as Priests. This capacity may have its deep roots in ancestor cults, but the predominant social role of Priests is found in politically integrated societies. Additional unpublished analyses implicate warfare in the decline of shamanism, and war powers implicate priestly roles, suggesting that warfare may be the real causal factor in the transformation of magico-religious practitioners. (3) Social Labeling and Malevolent Activities. If there was selection for the role on the basis of (a) negative social labeling alleging a biological inheritance of the role, or other forms of unwanted attribution then (b) the practitioner has an exclusively malevolent role characterization involving activities such as (c) causing illness, death and misfortune. These features were characteristic of the practitioners labeled Sorcerers/Witches. This negative dimension of the supernatural is produced through the persecutions carried out primarily by Priests and Healers, who designate people as being a Sorcerer/Witch. In your work, you often argue that religious practices and institutions (e.g., shamanic healing rituals) are deeply adaptive (Winkelman, 2009, Chapter 7 and 9, 2010a, Chapter 6; Winkelman \& Baker, 2010, Chapters 5, 6 and 11). Now, some cognitive scientists of religion have argued that religious practices and institutions consist mainly of spandrels and not of adaptations (e.g., Atran, 2002; Boyer, 2001; on the debate between adaptationism and spandrelism: Sosis, 2009). What are your main arguments in favor of the adaptationist account of religion and main objections against the spandrelist account? I have three chapters in an upcoming book The Super-Natural after the Neuroturn that I summarize in the following paragraphs. The simple evidence against the spandrel arguments is the many different functional and adaptive advantages of the diverse components of religious thought, as well as its group level effects. It is very easy to show the powerful social effects of supernatural assumptions, costly displays, ingroup cognition, etc. Similarly, the psychological benefits from enhanced endorphins, anti-depressive effects, enhanced social support, etc. speak to a lot of components of religious thought that facilitate human adaptation. Whether these are strictly biological or emerge in cultural-gene-environment interactions is really immaterial. The cognitive, social and physical environment is always part of the expression of genetic capabilities. So why consider religion adaptive? Michael Winkelman - The evolutionary neuroanthropology of consciousness 60 ALIUS Bulletin n3 (2019) aliusresearch.org/bulletin When primatologists study animal ritual they don't debate whether or not the rituals have adaptive functions. Why change the tune when we come to humans? Rituals exercise a variety of adaptive functions in moderating emotions, psychological and social relations and cognitive processes. Human's supernatural behaviors involve exaptations of the functions of primates' ritualized displays that were used to expand mechanisms for social communication and coordination as part of human innate psychosocial processes and cognitive structures. The communicative and integrative displays among primates, particularly the ritualized behaviors for group unification among great apes, provide a framework for understanding the origins of ritual behavior in activities that enhanced social integration. The concept of costly displays provides a framework for identifying the forces that led to the shamanic expansion out of the hominid ritual capacity, using drumming, singing and dancing to expand social integration function provided by mimesis. Shamanism emerged in this expansion of the mimetic capacity and its associated suite of expressive capacities that extended the social coordination functions of displays. Shamanic rituals expanded as adaptations involving increased capacities for ritual bonding of communities and the associated enhanced endorphin and placebo healing responses. The varied ways in which diverse alterations of consciousness contribute to healing attests to basic biological functions of this mode of consciousness. Shamanic alterations of consciousness reflect the physiological effects of ritual practices in stimulating the modulatory neurotransmitter systems of serotonin, dopamine and the endocannabinoids, as well as the endogenous opioid system. These provide the biological bases for these experiences in an enhancement of the functioning of neurotransmitter systems that enhance access to evolutionarily early strata of the brain. These brain areas provide the special cognitive qualities of consciousness that underlie perceptions of the supernatural. Ritual practices induce supernatural experiences through disrupting higher order information integration and top-down cognitive control, permitting emergence of cognitive processes related to ancient brain structures and primary process levels of cognition, identity and awareness. Shamanic alterations of consciousness provide adaptations that enhance cognition through expanded access to unconscious mental processes and the ability to integrate the global brain dynamics, rather than just the habitual networks (i.e., default mode network). The effects of psychedelics on global brain dynamics shows an enhanced global connectivity and increased connectivity between areas that are not normally connected. This illustrates how ASCs provide adaptive benefits by increasing access to novel information and cognitive processes. Michael Winkelman - The evolutionary neuroanthropology of consciousness 61 ALIUS Bulletin n3 (2019) aliusresearch.org/bulletin " The effects of psychedelics on global brain dynamics shows an enhanced global connectivity and increased connectivity between areas that are not normally connected. This illustrates how ASC provide adaptive benefits by increasing access to novel information and cognitive processes. It is not entirely clear to me whether you endorse a pluralist or unified account of altered states of consciousness. First, let me mention for the readers who are not familiar with your work that you distinguish between states of consciousness and modes of consciousness (Winkelman, 2010a, p. 127 et sq.). Your definition of modes of consciousness is close to that proposed by Tim Bayne and colleagues (Bayne \& Hohwy, 2016; Bayne, Hohwy, \& Owen, 2016)-your definition might be even broader (hyperonymic) than theirs. On the other hand, states of consciousness are to be understood as a more restricted hyponymic category. For example, according to you, soul flight and near-death-experiences are two distinct states of consciousness that belong to the same mode of consciousness, that you call the "integrative mode." Furthermore, in Shamanism, when you introduce the three main conscious types of states belonging to the integrative mode-the shamanistic type, the mystical type and the mediumnistic type-you define them as patterns (Winkelman, 2010a, p. 127). Therefore, it seems that what we have is a three-level taxonomy of consciousness. For a reconstruction of this three-level taxonomy, see Figure 1 (note that this tree is not exhaustive: several conscious modes, conscious patterns and conscious states are not mentioned). Figure 1: A reconstruction of the three-level taxonomic structure of consciousness proposed in Shamanism: A Biopsychosocial Paradigm of Consciousness and Healing everyday mode of consciousness integrative mode of consciousness shamanistic pattern of consciousness flight soul conscious state death and rebirth conscious state mystical pattern of consciousness mediumnistic pattern of consciousness " Michael Winkelman - The evolutionary neuroanthropology of consciousness 62 ALIUS Bulletin n3 (2019) aliusresearch.org/bulletin Now that these different concepts have been clarified, I can introduce the key idea of your neuroanthropological model of consciousness. According to you, no more than four modes of consciousness can be identified (Winkelman, 2010a, p. 22, 2011, p. 29): (1) Waking mode of consciousness. (2) Deep sleep mode of consciousness. (3) Dreaming mode of consciousness (REM sleep). (4) Integrative (spiritual and transpersonal) mode of consciousness (IMC, for short). A first question I have is: do you consider the dreaming mode as distinct from the integrative mode? Indeed, some passages of your work suggests the dreaming and the integrative modes may actually be closer than what your quadripartite taxonomy suggests. For example, you note that "[b]ecause of [...] similarities [between] REM and the IMC, shamans explicitly sought to integrate dream processes within ritual to induce alterations of consciousness" (Winkelman, 2010a, p. 136). Somewhere else, you say that the integrative mode is an extremely broad category including states as different as those induced by "hallucinogens, amphetamines, cocaine, marijuana, polypeptide opiates, long-distance running, hunger, thirst, sleep loss, auditory stimuli such as drumming and chanting, sensory deprivation, dream states, meditation, and a variety of psychophysiological imbalances or sensitivities resulting from injury, trauma, disease, or hereditarily transmitted nervous system conditions" (Winkelman, 2011, p. 31, my emphasis). Why classify "dream states" within the integrative mode of consciousness if, as previously stated, oniric states are purported to form an independent mode of consciousness? Why, in the first place, consider that the integrative mode (e.g., psychedelic experience or meditation) are different from dreaming? It is not clear to me why the taxonomic distance between dreaming and psychedelic experience is held to be greater than the distance between psychedelic experience and long-distance running, shamanic trance, meditation, hypnosis, etc. In fact, it seems to me that dreaming and psychedelic experience have much more in common (Kraehenmann, 2017; Kraehenmann et al., 2017) than psychedelic experience and amphetamine- induced arousal (Fink, 1969; Sanz, Zamberlan, Erowid, Erowid, \& Tagliazucchi, 2018). And yet, if I understand it correctly, your quadripartite taxonomy of conscious modes suggests otherwise! I think that the general physiological shift from extreme autonomous nervous system activation (either branch, especially sympathetic) to the rebound into a strong parasympathetic state provides a disengagement from the waking mode that enables access to many different structures of consciousness. So here a generalist model of ASCs; from there, however, intention, cognitive activity, exogenous influences (music, drums, drugs) can lead the person to a variety of states of consciousness. Michael Winkelman - The evolutionary neuroanthropology of consciousness 63 ALIUS Bulletin n3 (2019) aliusresearch.org/bulletin The literature on the persuasiveness of dreaming in mammalian species make it easy to see it is something separate from what I want to conceptualize in the integrative mode of consciousness. This is not to say that dream states are unrelated-indeed the whole concept of lucid dreaming provides a perspective on how a specific kind of awareness allows one to shift from the dream mode to an aspect of the integrative mode that includes self-awareness. So, integration of self-awareness (a waking mode capacity) into the dream mode may be sufficient to produce the integrative mode of consciousness, and extraordinary (but not ordinary) dream states may be one of the original platforms for selection for shamanic ASC. But clearly there are other ways to get into the ASCs of the integrative mode of consciousness, in particular by breaking down the coping mechanisms of the waking mode of consciousness-drugs, sleep deprivation, extreme pain, exhaustion, fasting, long-distance running-these all lead to a breakdown of the waking mode and allow for new discrete states of consciousness to emerge. I would not necessarily insist on only 4 modes, what is lumped within the integrative mode of consciousness may ultimately be shown to involve several different modes (i.e, sympathetic predominant versus parasympathetic predominant). I would see my last sentence here as my answer to your whole next question. Lead the way into more modes. But we have to be clear on what constitutes a mode of consciousness. I think a functional system is a good starting point. What can be done within a specific set of parameters of consciousness? Before going further, it might be useful to dwell a bit longer on your definition of what counts as a mode of consciousness. Your model draws its inspiration from Arnold Mandell's (1980) psychobiological model of consciousness. Mandell "suggested that physiological mechanisms underlying "transcendent states" are based in a common underlying neurobiochemical pathway" (Winkelman, 2010a, p. 25, my emphasis). Critical, then, is the idea that different phenomenological states may be underlain by a single neurobiological final pathway. The problem with this definition is that it remains rather vague as to how modes of consciousness should be parsed. In another passage, you point out that "modes of consciousness are revealed in the recurrent patterns of systemic neurophysiological functioning and their homologies with the major differences in experience" (Winkelman, 2010a, p. 22). So, to summarize, from what I understand, modes of consciousness are (i) underpinned by some basic organismic mechanisms-some basic "biological functions and organismic functions and needs" (Winkelman, 2011, p. 29)-and, (ii) these mechanisms give rise to various altered states of consciousness. Importantly, modes of consciousness are more encompassing than states of consciousness because modes refer to the underlying neurobiological pathways that have the potential of generating a large range of altered states (i.e., there is a one-to-many mapping between each mode and its characteristic states). We are now in a position to examine your proposal that there are exactly four modes of consciousness. More specifically, I would like to look at the third conscious mode: Michael Winkelman - The evolutionary neuroanthropology of consciousness 64 ALIUS Bulletin n3 (2019) aliusresearch.org/bulletin dreaming. Given what your definition of a conscious mode is, it could be argued that it would only make sense to distinguish between the non-dreaming deep sleep mode and the dreaming REM sleep mode if it were true that REM sleep and dreaming were perfectly overlapping with one another. However, as has been abundantly demonstrated within the last decades, a great deal of dreaming is in fact going on outside of REM sleep (Foulkes, 1962; Nir \& Tononi, 2010; Siclari et al., 2017; Stickgold, Malia, Fosse, Propper, \& Hobson, 2001). As a consequence, it seems that we are faced with a dilemma: (1) either we consider that dreaming as a whole constitutes a mode, but then the claim that there is a single neurobiological pathway underlying each mode will be violated because at least two pathways-REM sleep and non-REM sleep-will be recognized as generating dreaming; (2) or, we split dreaming in two-REM sleep and non-REM sleep-but then we are led to posit the existence in total of at least five modes of consciousness (the third one defined above being now split into two distinct modes). In sum, either the definition of what a mode of consciousness is should be revised or there are actually more modes of consciousness than previously recognized. Do you think the line of reasoning just sketched is sound? If so, which option would you be tempted to choose: revising the definition of modes of consciousness or recognizing the existence of more modes than initially proposed? Certainly there are more modes to be identified, and the integrative mode of consciousness may be too inclusive, requiring differentiation into biologically distinct functions of consciousness. It seems that the remarks just made about the dreaming mode of consciousness could also be made about the integrative mode of consciousness. In your work, you have identified many key phenomenological and neurobiological differences between various states belonging to the integrative mode. To cite a few examples: - Parasympathetic vs. sympathetic spiritual states: "While dominant meditative traditions emphasize the direct approach to parasympathetic dominant states, there are traditions that also engage the route of sympathetic stimulation, such as in the dancing of the Islamic mystics, the whirling dervishes." (Winkelman, 2010a, p. 131) - Meditation vs. shamanism: "[Meditators'] typical activities contrast with shamans in terms of more self-control and concentration, lower arousal, a sense of calm and emotional detachment, a loss of sense of self, a greater awareness, and contentless experience." (Winkelman, 2010a, p. 132) - Shamanism vs. possession: "Harner (1982) emphasized the shaman's remembering what happened during the soul journey as a characteristic of the shaman's SoC. Winkelman's (1986b, 1992) research supports the contention made by Eliade (1964) that this experience is not one in which the shaman is possessed by spirits, but, rather, one in which the shaman exercises a control over the spirits." (Winkelman, 2010a, p. 133) "A difference between shamanic flight and spirit possession is found in the association of the latter with amnesia. Although memory does occur in some situations of possession, amnesia does not occur with soul flight, and in the cross-cultural data, all Michael Winkelman - The evolutionary neuroanthropology of consciousness 65 ALIUS Bulletin n3 (2019) aliusresearch.org/bulletin cases of ASCs with amnesia are associated with possession." (Winkelman, 2010a, p. 173) All these differences seem to support a pluralist account of the integrative mode of consciousness. Indeed, the unity of this mode is arguably illusory and various "integrative" constructs should accordingly be identified. Thus, either more modes of consciousness should be recognized, or, alternatively, the same number of conscious modes could be retained, but then the definition of conscious modes should be drastically revised. What is your view on this dilemma? I would say that the question of the number of modes should be a question answered by empirical data. You have made the point about the need for a greater number of modes. Research by Fox et al. (2016) would support the expanded mode concept, even for meditation. Fox et al. found different patterns of brain activation and deactivation associated with different styles of meditation (focused attention, mantra recitation, open monitoring, and compassion/loving-kindness), as well as some similarities across most major meditation styles. Among their central findings was that the different categories of meditation had both unique psychological features as well as distinct patterns of activation and deactivation of different regions of the brain. A similarity across different forms of meditation involves effects on mental-physical processes that can be progressively developed through practice, engaging an ability to regulate physiological and mental activities, including involuntary processes. Different meditation techniques also recruited similar areas of the brain including: the insular cortex; the pre/supplementary motor cortices; the dorsal anterior cingulate cortex (involved in the regulation of attention and emotion); and the frontopolar cortex. Fox et al. reported that the primary meditation practices (except the loving-kindness/compassion technique) had effects on the posterior dorsolateral prefrontal, premotor and supplementary motor cortices. Fox et al.'s meta-analysis showed that practices of focused attention meditation produced significant activation both in prefrontal, premotor and dorsal anterior cingulate cortices, as well as slightly sub-threshold activation in the posterior dorsolateral prefrontal cortex and left mid insula. The focused attention meditation practices also resulted in a deactivation of two major default mode network hubs, the posterior cingulate cortex and the posterior inferior parietal lobule. Fox et al.'s meta-analysis of studies on open monitoring meditation techniques revealed significant activation in the insula, left inferior frontal gyrus, pre- supplementary and supplementary motor area, and premotor cortex, as well as posterior dorsolateral region of the prefrontal cortex and the dorsal anterior cingulate cortex. Michael Winkelman - The evolutionary neuroanthropology of consciousness 66 ALIUS Bulletin n3 (2019) aliusresearch.org/bulletin So this kind of specific brain system differences is the kind of evidence that would support breaking down into further modes of consciousness. But this would have to be on functional grounds too that I discuss below. The reason why you call the integrative mode of consciousness integrative is that this mode is characterized by the integration of distinct informational units of the brain that do not usually communicate with one another (or do so to a much lesser degree). In your own words, this conscious mode "produce[s] an integration of information processing between the R-complex [the reptilian brain] and the limbic system, between the limbic system and the frontal cortex, and between the hemispheres of the cortex" (Winkelman, 2011, p. 38). Overlapping with the concept of integration is the concept of psychointegrator. The latter refers to any chemical compound having the propensity to cause integration in the brain and in particular to "provoke limbic discharge patterns that produce enhanced interhemispheric synchronization and increased communicative interaction between frontal hemispheres, and between the lower brain areas and frontal cortex" (Winkelman, 2001, p. 220). In other words, "psychointegrator" is a less pejorative and more technical way to speak of "serotonergic hallucinogens" (LSD, mescaline, psilocybin, DMT, etc.) (Winkelman, 1996, 2001, 2007). Now, what is the evidence in favor of the view that serotonergic hallucinogens have a psychointegrative effect? Let us look first at the EEG data. One central characteristic of the integrative mode of consciousness is that it is dominated by highly synchronized and coherent brain waves. In particular, we would expect the EEG signal to be dominated by the high-frequency gamma and low-frequency theta rhythms (e.g., Winkelman, 2010a, p. 35, 2011, p. 30). Do psychointegrators effectively induce such states? The answer seems to be no. There is now a plethora of EEG studies available on serotonergic hallucinogens, but none of them describe the gamma rhythm as being particularly increased by the intake of a psychedelic compound. To my knowledge, the only exception to the rule is Eduardo Schenberg et al.'s (2015) EEG study of ayahuasca. However, this increase in gamma was restricted to the second half of the experience, and importantly, it was mainly interpreted by the authors as being caused by harmaline (Schenberg et al., 2015, pp. 20-21), which is not a serotonergic hallucinogen (i.e., a psychointegrator) and whose activity probably stems chiefly from anticholinesterasic mechanisms (Yang et al., 2015; Zhao et al., 2013; Zheng et al., 2009). As regards the increase in slow rhythms-and theta in particular-again, to my knowledge, this electrophysiological change does not typically characterize psychointegrators (e.g., Muthukumaraswamy et al., 2013); instead, it is typical of antimuscarinic hallucinogens (Ebert, Grossmann, Oertel, Gramatte, \& Kirch, 2001; Itil, 1966; Itil \& Fink, 1966; Osipova et al., 2003), whose phenomenology and neurophysiology is completely different from that of serotonergic hallucinogens (Fortier, 2018a, 2019; Gyermek, 1998; Ketchum, Sidell, Crowell, Aghajanian, \& Hayes, 1973). So, looking at the EEG data, it could be objected that typical "integrative rhythms" are not present in psychointegrator-induced states. What is your take on this matter? Michael Winkelman - The evolutionary neuroanthropology of consciousness 67 ALIUS Bulletin n3 (2019) aliusresearch.org/bulletin The measurement of brain waves and their various forms of coordination and coherence is beyond my expertise. I think that nonetheless, the idea that ASC in general, and psychedelics included, reduce the overall brain wave frequencies is supported by diverse forms of evidence. I note you refer to "plethora of EEG studies available on serotonergic hallucinogens" but I did not see references. Here, let me quote at length some already published material where I directly address these issues (Winkelman, 2017a, 2017b, 2018a): In spite of the diversity of entheogenic species and the broad range of psychoactive substances, the principal psychedelics share similarities as tryptamines and indole alkaloids, which are sources of DMT and similar neurochemicals that function as agonists stimulating the serotonergic system. Serotonin has been considered the primary neurotransmitter system affected by psychedelics, especially through their effects at 5-HT2A receptors; action on other serotonin (5-HT) receptors is also established, as well as a wide range of other neurotransmitter systems. The phasic effects of psychedelics first stimulate and enhance serotonin; secondly, saturate and overload the serotonin system; and thirdly, release the habitual serotonin repression of the dopaminergic system. Psychedelics' resistance to normal reuptake mechanisms locks out serotonergic transmitter sites, habituating the receptors and reducing the regulatory processes of the serotonergic system. This results in a release of the dopamine system normally repressed by serotonin, causing a variety of visionary experiences (hallucinations, dreams, psychosis) and modifying control and coordination among the major brain subsystems. Psychedelics compromise the serotonergic inhibition of the ascending flow of information and emotional responses, resulting in the release of information from ancient levels of the brain that is normally inhibited by serotonin. These effects are typified by psychedelics' interruption of cortico-striato-thalamo-cortical loops that inhibit the lower brain structures' sensory gating systems, providing an enhanced availability of information managed by these brain areas (Vollenweider, 1998; Vollenweider \& Geyer, 2001). These psychedelic effects in altering consciousness are illustrated by Vollenweider's (1998) research on the mechanisms of action of psychedelics on the major cortical loops. The frontal-subcortical circuits provide one of the principal organizational networks of the brain involving neuronal linkages and feedback loops of the cortical areas of the frontal brain with the thalamus of the brain stem region. Vollenweider's attributes the consciousness-altering properties of psychedelics to their selective effects on Michael Winkelman - The evolutionary neuroanthropology of consciousness 68 ALIUS Bulletin n3 (2019) aliusresearch.org/bulletin the brain's cortico-striato-thalamo-cortical feedback loops that link the information gating systems of lower levels of the brain with the frontal cortex. The typical action of psychedelics interrupt the cortico-striato- thalamo-cortical loops that inhibit the lower brain structures' sensory gating systems that reduce the flow of information to the frontal areas of the brain (Vollenweider \& Geyer, 2001). Psychedelic interruption of serotonergic inhibition of thalamic screening results in a flood of information from these ancient levels of the brain. This overwhelms the processing capacities of the frontal cortex and leads to alteration of experience of self, other, environment and produces a focus on the internal world of psychological structures and projections. The inhibition of dopamine release by serotonin is central to neurochemical balance in the brain, with the serotonergic and noradrenergic systems of the right hemisphere inhibiting the dopamine system and the left hemisphere (Previc, 2009). This blockage of serotonin's inhibitory functions results in the disinhibition of the dopaminergic system, releasing a flood of information that is normally inhibited by serotonin. The reduction of serotonergic and noradrenergic modulation (control) results in the ascendance of the dopaminergic and acetylcholine systems that produce a variety of notable visual syndromes, especially hallucinations and dreaming (Hobson, 2001). Psychedelics and other alteration of consciousness share common effects on the brain's perceptual mechanisms and representational capacities through a temporary deregulation of the prefrontal cortex (PFC). Common effects of this disruption are manifested in the loss of the roles of the frontal lobes and prefrontal cortex involving higher cognitive functions. This disruption of the PFC results in the loss of various capacities-capacity for willful action, deliberate direction of attention, and aspects of self-awareness, as well as the capacities for abstract thought, creativity and planning. When these higher level brain functions of the cortical regions and the PFC are down regulated, this allows for the manifestation of lower brain structures usually repressed by the PFC. This means the emergence of information and aspects of identity that are related to our more ancient brain regions. Psychedelic disruption of the DMN permits the operation of a more fluid and dynamic brain lacking its usual top-down principles of control. Psychedelics such as LSD, psilocybin and ayahuasca cause decreases in DMN brain activity (Carhart-Harris et al., 2012; Carhart-Harris et al., 2014; Carhart-Harris et al., 2016; Palhano-Fontes et al., 2015) and the disintegration of normal DMN functions. This is the consequence of a reduction in the connectivity of the frontal cortex with lower brain areas (Alonso, Romero, Mananas, \& Riba, Michael Winkelman - The evolutionary neuroanthropology of consciousness 69 ALIUS Bulletin n3 (2019) aliusresearch.org/bulletin 2015) and the reduction in oscillatory activity and power in posterior and frontal association cortices (Muthukumaraswamy et al., 2013). This involves a decrease in the functional coupling of the frontal cortex with the medial temporal lobe, as well as of the medial prefrontal cortex with the posterior cingulate cortex. Consequently, the lower brain dynamics involving ascending circuitry are released, providing a strong input to the frontal cortex from the ancient brain systems. This dynamic is hypothesized as the mechanism that releases the innate modules and promotes their manifestations in consciousness (also see Winkelman, 2017b). A principal effect of psychedelics involves production of hypersynchronous ascending slow wave brain discharges in the hippocampal-septal-reticular- raphe circuit that impose impulses from the ancient lower stratum of the brain on the frontal areas (Mandell, 1980). This pattern of psychedelic action on the brain is shared by many other agents and conditions that alter consciousness (Winkelman, 2011). Alterations of consciousness produced by behavioral and physiological conditions, as well as pharmacological agents, causes a reduction in the serotonin inhibition to the hippocampal cells, which results in an increase in slow-wave EEG activity in the hippocampal-septal region. Research on the mechanisms of action of psychedelics on the major cortical loops (Vollenweider \& Geyer, 2001) illustrates these effects. Psychedelic effects on the cortico-striato-thalamo-cortical (CSTC) feedback loops and their regulatory effects on the gating systems of lower levels of the brain alter consciousness. This psychedelic interruption of the CSTC loops and their inhibitory function release the lower brain structures' sensory gating and enhance the flow of information to the frontal areas of the brain. Psychedelics interruption of thalamic screening results in a flood of information from these ancient levels of the brain. Alonso et al. (2015) confirmed this psychedelic enhancement of a bottom-up information dynamic by psychedelics with ayahuasca. They assessed ayahuasca-induced changes in directionality of information flow in the brain, with changes in the connectivity of brain oscillations. These changes involved a disruption of the normal coupling between anterior and posterior areas of the brain that resulted from reduction in the influence of frontal brain areas over the posterior areas. This reduction was accompanied by increases in the influence of posterior brain regions on the frontal anterior areas. "These results suggest that psychedelics induce a temporary disruption of neural hierarchies by reducing top-down control and increasing bottom-up information transfer in the human brain" (Alonso et al. 2015, p. 1). Michael Winkelman - The evolutionary neuroanthropology of consciousness 70 ALIUS Bulletin n3 (2019) aliusresearch.org/bulletin The release of these lower brain areas is the likely cause of the enhanced operation of the innate modular operators of the brain. These operators reflect unconscious cognitive processes that were acquired deep in evolution, with some apparently widely shared by other primates and mammals. These ancient roots of these operators suggest that they are associated with the function of our ancient brain structures. Winkelman (2010a) has detailed how many features of shamanic alterations of consciousness can be explained by reference to the operation of these ancient brain structures, particularly the paleomammalian brain. Regardless of the question of knowing which electrophysiological rhythm underlies the integrative mode of consciousness one might be tempted to challenge the view that such a mode is indeed integrative. In your initial definition, what you meant by integration was that different parts of the brain that are usually disconnected suddenly become connected. An obvious way to measure such (dis)connectivity in the brain is to look at functional connectivity (Friston, 2011). Fortunately, in recent years, several studies have looked at this dimension of the psychedelic experience. By and large, the pattern of finding is that functional connectivity is locally decreased but globally increased, a point that you duly acknowledge and nicely discuss in your recent work (Winkelman, 2017b, pp. 7-8). For example, it has been shown that within the default mode network, connectivity is weakened, but conversely different areas of this network start talking to other areas and networks from which they are usually disconnected (Carhart-Harris et al., 2014). Moreover, in some cases, the large-scale connectivity between networks can even be decreased: e.g., with psilocybin, the connectivity between visual and sensorimotor areas is weakened (Roseman, Leech, Feilding, Nutt, \& Carhart-Harris, 2014). The key lesson is that overall psychointegrators do not maximize connectivity. At this point of the discussion, it will be useful to introduce two concepts: complexity and neural criticality. The first concept was developed by Giulio Tononi (Tononi, Edelman, \& Sporns, 1998; Tononi, Sporns, \& Edelman, 1996). It is a measure combining two indexes: segregation and integration. If areas of the brain are excessively segregated, complexity will be low and this will result in a lack of conscious experience (Boly et al., 2012; Tagliazucchi et al., 2013) (see bottom right circle in Figure 2). On the other hand, if areas of the brain turn out to be excessively integrated, complexity will also be too low for consciousness to emerge (Arthuis et al., 2009; Blumenfeld, 2012) (see bottom left circle in Figure 2). Complexity is the highest when segregation and integration are optimally balanced. This is typically the case in the everyday waking state. Michael Winkelman - The evolutionary neuroanthropology of consciousness 71 ALIUS Bulletin n3 (2019) aliusresearch.org/bulletin Figure 2. Three combinations of brain segregation and integration (from: Tononi et al., 1998, p. 478) Another concept, criticality, has many points in common with that of complexity (Tagliazucchi \& Chialvo, 2013; Timme et al., 2016). It refers to power-law distributed physical phenomena that often take place in phase transition. The study of neuronal avalanches in the brain has also revealed the presence of such power-law distributed firing of neurons (Beggs \& Plenz, 2003; Chialvo, Balenzuela, \& Fraiman, 2008). Interestingly, critical states largely overlap with complex states; furthermore, Tononi's excessively segregated states share many properties of sub-critical states and excessively integrated states with super-critical states. Now, recent research conducted mainly at Imperial College has demonstrated that psychedelic states are very close to criticality-in fact, even closer than normal waking states (Atasoy et al., 2017; Carhart-Harris, 2018; Carhart-Harris et al., 2014). This is consistent with the aforementioned findings to the effect that serotonergic hallucinogens do not increase connectivity simpliciter but rather do so in a very balanced way (connectivity is rearranged through both increase and decrease of connectivity patterns). It appears, then, that psychedelics are not properly speaking psychointegrators. A true psychointegrator would be a compound able to trigger states with the highest integration and the lowest segregation possible (bottom left circle in Figure 2). It seems that this is not what serotonergic hallucinogens do; instead, this is what epileptic seizures do (Arthuis et al., 2009; Meisel, Storch, Hallmeyer-Elgner, Bullmore, \& Gross, 2012; Priesemann, Valderrama, Wibral, \& Le Van Quyen, 2013). Therefore, don't you think that only epileptic seizures should deserve the title of "integrative mode of consciousness" (i.e., "supercritical mode of consciousness")? And don't you think that serotonergic hallucinogens should rather be called "psychocomplexifiers" or "psychocriticalizers" and the mode of consciousness they induce the "complex or critical mode of consciousness?" Michael Winkelman - The evolutionary neuroanthropology of consciousness 72 ALIUS Bulletin n3 (2019) aliusresearch.org/bulletin Well any integration implies a deafferentiation or separation from something else. So what is being integrated? Good evidence that a shared functional feature of ASC is the enhanced integration of lower brain structures and activity, the ascending stimulation of the cortico-striato-thalamo-cortical loops and the serotonergic hippocampal-septal system. So that is what is being integrated, lower brain processes into the frontal brain. There are certainly other aspects of global integration that occur as well that you mentioned in the first part above, the globally increased functional connectivity and increased connectivity of areas and networks from which they are usually disconnected. So I would say this new generation of research is supporting the notion of some forms of enhanced connectivity, albeit at the expense of breakdown of the normal connectivity networks of the default mode network and prefrontal cortex. There is arguably another reason for thinking that psychedelics are not genuine psychointegrators. According to you, the integrative mode of consciousness coincides with a "loss of hippocampal CA3 modulation removes regulatory input from the environment, leaving the 'inside world' dominant" (Winkelman, 2010a, p. 26, my emphasis). In another passage you also argue that "[t]he highly integrated internal visionary states of the shaman involve a reduction of other inputs, in essence a reduction of external sensory information and bodily input" (Winkelman, 2010a, p. 32, my emphasis). In sum, in the integrative mode of consciousness, the central nervous system becomes less sensitive to the "external world" and more sensitive to the "internal world" (i.e., to itself). If so, then, it does not seem that serotonergic hallucinogens are psychointegrators, because like any (near) critical state, they induce an increased sensitivity to stimuli from the external world (Atasoy et al., 2017; Carhart-Harris, 2018). As Enzo Tagliazucchi puts it, "physical systems at criticality present a maximal susceptibility, i.e. a maximal response to external perturbations" (2017, p. 145). In other words, it appears that psychedelics do the contrary to what a psychointegrator is purported to do: they increase-rather than decrease-responses to the external world. Don't you think this supports the view that psychointegrators (a.k.a. psychedelics) would be better described as psychocriticalizers? Since the earliest studies of the psychedelics in the context of psychiatry it was recognized that the primary determinants of the effects are from set and setting, the individual's expectations and the influences in the environment. So there is a lot of variation in terms of what psychedelics can produce. Some cultures used ayahuasca as preparation for headhunting and warfare, dancing for hours before leaving during the night to carry out raids. Other cultures use ayahuasca for journeying across the cosmos or healing. The same substance can produce diverse outcomes, depending on how you engage with it. There is also some confusion promoted in recent brain research from a failure to pay attention to the mode of administration of Michael Winkelman - The evolutionary neuroanthropology of consciousness 73 ALIUS Bulletin n3 (2019) aliusresearch.org/bulletin psilocybin-whether it is drunk or injected. It makes a big difference whether you chew coca leaves or inject cocaine, we should expect the same for psilocybin. You have endorsed a strong perennialist approach in your work on the contents of the integrative mode of consciousness. For example, according to you, mystical and psychedelic states are very similar because they both tap into the same evolutionary mechanisms: "The relationships among natural and drug-induced alterations of consciousness must be understood from an evolutionary perspective. This reveals altered consciousness to be related to endogenous mechanisms that are triggered by both ancient evolutionary adaptations and more recently acquired propensities to use exogenous sources of substances to alter consciousness" (Winkelman, 2010a, p. 27). All integrative states seem to share common themes, images, and features. As you note, "[t]he apparent similarities in psychedelic entities and various other types of entity experiences found across cultures, time, and diverse conditions for altering consciousness suggest that an explanation be sought within innate functions of the human brain" (Winkelman, 2018a, p. 3). Such a view echoes the Jungian theory of archetypes according to which the cognitive and perceptual features that are universally found across cultures are to be explained in terms of underlying unconscious archetypes (Laughlin, 2011, Chapter 10; Laughlin et al., 1990, p. 134; Winkelman, 2010a, p. 220, 2018a, p. 7; Winkelman \& Baker, 2010, pp. 194-198). Your perennialist evolutionary approach raises two questions: (1) why should universal cognitive features be necessarily explained in evolutionary terms?; and, (2) what is the actual evidence in favor of perennialism? Let me start by examining the first question. Is it true that universality is a cue to innateness? Well, if it is, it could be argued that it is a very inaccurate cue. There are plenty of cognitive and perceptual phenomena that are universal and probably not innate. For example, all humans have the prior expectation encoded in their visual system that light is coming from above (Sun \& Perona, 1998). Among other things, the famous concave/convex circle illusion is explained by this prior. Given the pervasiveness of such a prior one may expect it to be innate. But it seems that it is not since it can be changed by experience (Adams, Graf, \& Ernst, 2004). The universality of the light-from-above prior appears to be learned. However, given that the structure of the environment is everywhere the same (i.e., light always comes from above in non- artificial environments), in this case, learning (absence of innateness) coincides with universality. As students of natural scene statistics have amply demonstrated, many universal features of the perceptual system and of perceptual experience can be explained by non-innate learning mechanisms combined with stable environmental structures (Geisler \& Kersten, 2002; Geisler, 2008). Another problem with the view that universality can be used as a cue to innateness is that it overlooks the possibility of diffusion phenomena. For instance, many apparently universal features of myths can be thoroughly explained by cultural diffusion. To take a specific example, some authors have proposed evolutionary accounts of the widespread dragon motif-both in iconography and myths and legends (Jones, 2016; Sperber, 1996b; Wengrow, 2013). If this motif is so pervasive, it is argued, it is because our brain is structured in such a way that it easily triggers and encodes chimeric Michael Winkelman - The evolutionary neuroanthropology of consciousness 74 ALIUS Bulletin n3 (2019) aliusresearch.org/bulletin creatures. But such an explanation sounds way too ad hoc. As a matter of fact, it has been shown that the distribution of the mythological dragon motif can be accounted for merely in diffusionist terms (d'Huy, 2013, 2014, 2016). The phylogeny of cultural representations reveals how the dragon motif has been gradually modified and updated from its earliest versions (in African cultures) to its latest versions (in Amerindian cultures). More generally, as demonstrated by population genetics, the phylogeny of myths seems to largely overlap with prehistoric human migrations (Korotayev \& Khaltourina, 2011). This overlap provides strong evidence in favor of the pervasiveness of diffusion phenomena. As a result, many universal traits of human culture can be explained by resorting to the prehistory of migrations and borrowings rather than by positing cognitive adaptive functions. In sum, it seems that the burden of proof is on the proponent of psycho-evolutionary and nativist explanations. For example, when you write that "[t]he widespread manifestation of these visual elements under diverse circumstances attests to the elicitation of these innate mechanisms in producing these characteristic psychedelic experiences" (Winkelman, 2017b, p. 9), it could be objected that pinpointing the universal distribution of a trait is not enough to demonstrate its evolutionary origin. In fact, some evolutionary psychologists have even called for a thorough dissociation between being an innate adaptive function and being universal (e.g., Apicella \& Barrett, 2016; C. Barrett \& Kurzban, 2006). What is your view on this issue? "Why should universal cognitive features be necessarily explained in evolutionary terms?" Well if you really have a more compelling line of explanation besides the innate tendencies and their adaptive potentials, then offer it. I think the combination of the cross-cultural manifestations in religion and their various putatively adaptive functions is the most compelling argument. What would you propose? Let me now turn to the second question: what is the actual evidence for perennialism? First, it should be said that Jung, one of the main inspirations of neuroanthropological perennialism, based his theory on a somewhat biased ethnographic database: he misinterpreted several ethnographic cases and carefully ignored numerous counter- examples which refute the theory of archetypes (Le Quellec, 2013, Chapters 8-9). But rather than focusing on the details of Jung's theory it is more interesting to examine what the evidence for perennialism is in the specific case of hallucinogenic experiences. According to you, hallucinogens offer a pharmacological model of mysticism whose relevance goes far beyond drug-induced mystical states: Similarities in psychedelic-induced visionary experiences and those produced by practices such as meditation and hypnosis and pathological conditions such as epilepsy indicate the need for a general model explaining visionary experiences. (Winkelman, 2017b, p. 1) The fundamental similarities of psychedelic-induced and naturally-induced mystical experiences (Smith, 2000; Yaden et al., 2017) support the classic Michael Winkelman - The evolutionary neuroanthropology of consciousness 75 ALIUS Bulletin n3 (2019) aliusresearch.org/bulletin perennialist view of fundamental commonalities to mystical experiences across cultures and their independence of the mode of induction. (Winkelman, 2017b, p. 5) It seems that the conflation of hallucinogenic states and other states like epileptic seizures and meditation could be resisted on several grounds. First, as we have seen before, from a brain-system perspective, epileptic seizures are very different from serotonergic-induced states: the former are supercritical whereas the latter are (near) critical. By the same token, although meditation and serotonergic-induced states have certainly many things in common, they nevertheless exhibit striking neurophysiological differences (Milliere, Carhart-Harris, Roseman, Trautwein, \& Berkovich-Ohana, 2018). Even without looking at the differences between hallucinogenic and non-hallucinogenic altered states and focusing only on hallucinogenic experiences, perennialism does not seem very compelling. Namely, when looking at the diversity of hallucinogenic experiences, it is difficult to see what kind of phenomenological feature truly fit with the theory of perennialism. Before further elaborating on this anti-perennialist argument, it will be helpful to clarify what is meant here by "hallucinogens." In your article on psychointegrators you put forward the following definition: Classification as a hallucinogen (or psychedelic) has not been based upon specific chemistry or physiology, but upon effects on human experience, producing visions, voices and effects upon perception, mood, and thought in non- toxic doses (Siegel, 1984). This distinguishes them from substances that produce hallucinations because of toxicity. (Winkelman, 2001, p. 220, my emphasis) Thus, according to this definition, hallucinogens and psychedelics are synonymous. This is no surprise given that you cite serotonergic hallucinogens (e.g., LSD, mescaline, psilocybin) as prototypical examples of hallucinogens (a.k.a. psychedelics). What is less clear is why you suggest hallucinogens should be restricted to compounds whose "mechanisms of action involv[e] intervention in serotonin pathways" (Winkelman, 2001, p. 220). Why not also include other classes of hallucinogens whose mechanisms are not serotonergic? Following the above definition of hallucinogens, it seems that other neuropharmacological classes should also be included: notably, -opioid hallucinogens (e.g., salvia divinorum), antiglutamatergic hallucinogens (e.g., ketamine), antihistaminergic hallucinogens (e.g., diphenhydramine), antimuscarinic hallucinogens (e.g., scopolamine), and hybrid hallucinogens (e.g., ibogaine). As we can see, there are various classes of hallucinogens and all of them fulfill the definitional criteria that you put forward in the above passage-including, it must be stressed, antimuscarinic hallucinogens, which can indeed be toxic at high doses (Winkelman, 2001, p. 232, f.n. 1), but which remain quite safe at hallucinogenic doses: as a matter of fact, the hallucinogenic threshold of these compounds is much lower than the toxicity threshold (Gyermek, 1998, pp. 352-353). For example, in the experiments at Edgewood Arsenal, the highest doses of BZ (3-Quinuclidinyl benzilate) given to participants were never higher than 10\% of the LD50 (the median lethal dose). Most of the time, doses that were administered did not exceed 3\% of the LD50 and yet copious hallucinations were already Michael Winkelman - The evolutionary neuroanthropology of consciousness 76 ALIUS Bulletin n3 (2019) aliusresearch.org/bulletin observed at such doses (Panel on Anticholinesterase Chemicals, Panel on Anticholinergic Chemicals, Committee on Toxicology, Board on Toxicology and Environmental Health Hazards, \& Assembly of Life Sciences, 1982, p. 61). Now, the objection that could be raised against any perennialist account of hallucinogens would go like this: you argue that the phenomenology of hallucinogens support perennialism (the view that all hallucinogenic experiences-and other non- chemically-induced altered states-share a common phenomenological core), but this is because you take only serotonergic hallucinogens into account; if you were to include in your analysis other classes of hallucinogens, you would then realize that massive differences exist between these compounds, and as a result, you would acknowledge that perennialism does not hold water. Let me flesh out this anti-perennialist objection with three concrete examples. It is very common for people to find themselves transformed into artefacts under the effect of salvia divinorum (they are typically metamorphosed into wheels or pieces of furniture surrounding them); by contrast, this is never reported in serotonergic-induced trips (Fortier, In preparation, Chapter 4). Similarly, antimuscarinics strikingly differ from serotonergics in their effects. While the latter often induce extraordinary hallucinations (spirits, anthropomorphic entities, chimeras, etc.) the former only induce hallucinations of very ordinary and mundane objects (humans, animals, everyday artefacts, etc.) (Fortier, In preparation, Chapter 4, 2018a, 2018b, 2019; Ketchum, 2006; Ketchum et al., 1973). Even entoptic hallucinations (Billock \& Tsou, 2012; Kluver, 1966; Lewis-Williams \& Dowson, 1988), which are often taken to be a universal feature of hallucinogenic experiences, are completely absent from antihistaminergic- and antimuscarinic-induced hallucinations (Fortier, In preparation, Chapter 4). After having examined the outstanding diversity of features induced by each class of hallucinogens, it is very difficult to see what kind of phenomenological feature could be argued to constitute the common perennial core of these experiences. Importantly, if this criticism of perennialism is correct, then hallucinogens (as defined above) cannot be said to form a single mode of consciousness anymore. Moreover, the contents of these hallucinogenic experiences cannot be said to share a common core with other nonchemically-induced mystical states because what is meant by "hallucinogenic experience" covers very heterogeneous experiences with no identifiable common core. Is this line of criticism of perennialism sound to you? If so, do you think that the evolutionary perennialist account of hallucinogenic experiences can nonetheless be salvaged from this anti-perennialist objection? "What is the actual evidence in favor of perennialism?" The universal features of religious and spiritual belief and practice, the cross-cultural variation in such patterns as a function of social complexity, the cross-cultural features of shamans, the universal manifestations of shamanistic healers, the fundamental similarity in many forms of meditative experience across diverse and disparate traditions, the Michael Winkelman - The evolutionary neuroanthropology of consciousness 77 ALIUS Bulletin n3 (2019) aliusresearch.org/bulletin similarity in mystical phenomenology across cultures and time, similarities in the structures and functions of ASCs induced by diverse mechanisms, etc. Here again, let me quote some published material where I directly address these issues (Winkelman, 2017a, 2017b, 2018a): 5-HT2A has been considered the primary neurotransmitter system affected by psychedelics such as LSD but action on other serotonin receptors and other receptor systems has also been established (Halberstadt \& Geyer, 2011; Ray, 2010, 2013, 2016). This is not to say that psychedelics only share common effects, they also have distinctive effects, neurologically as well as phenomenologically (see: Ray, 2010, 2016). In assessing the different profiles of neurotransmitter interactions by various types of psychedelics, Ray (2016, p. 49) noted that "[m]ost of the drugs studied interact with multiple receptors, and most of the receptors studied interact with multiple drugs" This leads Ray (2010, 2012) to challenge the dominant theory of psychedelic action as being primarily mediated by effects at the 5-HT2 receptors. Examination of the relative affinity of various psychedelics for a wide range of receptors (Ray 2010, p. 22 and 41) found that "LSD has the strongest interaction collectively with the five dopamine receptors [...] [and] DMT has the strongest interaction with any single dopamine receptor. Rolland and colleagues (2014) noted that hallucinations-or perhaps less pejoratively "visionary experiences"-may be induced by a variety pharmacological mechanisms, including the hyperactivation of dopamine receptors, such as that caused by psychostimulants; stimulation of serotonin 5HT2A receptors targeted by psychedelics; and the blockage of glutamate NMDA receptors caused by dissociative anesthetics. They proposed that these different pharmacological systems might share common neurobiological pathways involving integrated neurobiological circuits that when compromised can produce hallucinations. And they hypothesized that diverse mechanisms, including dopamine and serotonin activation and NMDAR blockage, can disrupt the thalamic gating functions and cortico- striato-thalamo-cortical loops, resulting in a disorganization of the brains basic filtering processes, and consequently leading to visionary experiences. Previc (2009) also proposes that diverse visionary alterations of consciousness are a function of the dopamine system, which is directly stimulated by many different neurotransmitters and drugs and indirectly through effects on other neurotransmitter systems. Previc (2006, 2009) reviews evidence indicating a common underlying mechanisms for diverse methods of altering consciousness involving a disinhibition of dopaminergic extrapersonal brain Michael Winkelman - The evolutionary neuroanthropology of consciousness 78 ALIUS Bulletin n3 (2019) aliusresearch.org/bulletin systems, particularly those in the ventral cortex and the limbic circuit. Independent of the specific neurotransmitters involved, diverse processes producing visionary experiences may share common underlying mechanisms in a thalamic sensory overload, the common pathway resulting in a disruption of cortico-subcortical processing (Vollenweider, 2001). This common pathway can carry a variety of different features of experience produced by the distinctive qualities of the various psychedelic substances. Ray (2012, 2016) proposes that the diversity in the phenomenology of psychedelic experiences is a consequence of the distinctly different neurotransmitter receptor profiles that each substance engages. Each neurotransmitter system (i.e., the various serotonin receptor subtypes, beta receptors, dopamine, histamine-1, imidazoline-1, kappa, mu, sigma, and cannabinoid receptors) elicits a specific profile of effects that Ray calls a mental organ. Each psychedelic drug has effects on a range of neurotransmitter systems that results in its characteristic effects on neuronal activity and on consciousness. His idea is akin to innate modules and modes of consciousness, but with a highly individualistic twist-each neurotransmitter is a distinct state at least, but a mode of consciousness? I would doubt that every neurotransmitter or receptor would each constitute separate modes, but they could. But the fact that each transmitter can affect different kinds of neurotransmitter systems undermines that possibility of a pure serotonin effect, for instance. I do not think that the individual neurotransmitters are really the best level of analysis for differences in consciousness. There are major global contrasts with the waking mode of consciousness that allows us to provide a generalized profile that can be induced by diverse mechanisms. It is these generalities that characterize ASC-parasympathetic dominance, internal focus of attention, visual experience, emotional and egoic activation, right hemisphere dominance, etc. that provide commonalities for diverse states of the integrative mode of consciousness. If you enter into a parasympathetic visionary state, the functional dynamics of that state are generally the same whether or not you got there by LSD or ibogaine or ayahuasca or some dopaminergic or GABA action. " I do not think that the individual neurotransmitters are really the best level of analysis for differences in consciousness. There are major global contrasts with the waking mode of consciousness that allows Michael Winkelman - The evolutionary neuroanthropology of consciousness 79 ALIUS Bulletin n3 (2019) aliusresearch.org/bulletin us to provide a generalized profile that can be induced by diverse mechanisms. So this gets back to how many modes do we need. Perhaps the answer is in terms of access. Most people can't access the dreaming mode directly from waking consciousness and vice versa, although some dream events do produce a rapid transition or warp between modes-as well as states. But you can't easily go directly from one mode to another unless you have special training, like in lucid dreaming. So how to apply this to the differentiation of modes and meditative states? How many modes do we need to explain meditative experiences? The first criteria ought to be access. Can you go willingly from void to bliss? If so, then to me the notion of different states of consciousness in the integrative mode of consciousness holds. If you can't go from one state of consciousness to another and they are functionally different in terms of what they can do, then different modes are involved. I suspect that much of this has been worked out in various Hindu and Buddhist traditions. I must confess my original notion of modes of consciousness came from some unremembered scholar of meditation. Certainly many Hindi scholars have proposed concepts similar to my modes of consciousness in the recognition of deep sleep, dream, waking and one or more transcendental forms of consciousness. These scholars may well have worked out this question of modes for us on the basis of phenomenology and function. Much of your research conducted in the last two decades was dedicated to the study of hallucinogens. While it is very common in the field of psychedelic studies to endorse a non-naturalistic framework positing that hallucinogenic entities have some kind of objective reality (e.g., Luke, 2017; Strassman, Wojtowicz, Luna, \& Frecska, 2008), you have played a key role in defending and promoting a naturalistic account of these experiences (other accounts following this path notably include: Kent, 2010; Letheby, 2016, Chapter 4). In a recent paper, you have very explicitly advocated this naturalistic approach: [...] we need an academic study of entity encounters that offers a thorough examination of the similarities in independent reports by identifying the recurrent characteristics common to these experiences. (Winkelman, 2018a, p. 6) [...] if we simply accept the phenomenological experiences of entities as transcendent realities, we commit an error of epistemological naivety. (Winkelman, 2018a, p. 5) " Michael Winkelman - The evolutionary neuroanthropology of consciousness 80 ALIUS Bulletin n3 (2019) aliusresearch.org/bulletin In your view, what are the key requirements that any proper naturalistic and scientific account of psychedelic experiences should satisfy? And what are the pitfalls that it should thoroughly avoid? Requirements for a mature psychedelic science must grapple with on one hand, a recognition of a unitary basis for all kinds of knowing and believing, while on the other hand recognizing what is the distinctive manner of knowing opened by psychedelics. So we need a perspective that combines unity and diversity in knowledge, and I think that means a necessary engagement with metaphysical dualism. What we need to avoid is ontological and metaphysical naivety. J " Requirements for a mature psychedelic science must grapple with, on one hand, a recognition of a unitary basis for all kinds of knowing and believing, while on the other hand recognizing what is the distinctive manner of knowing opened by psychedelics. [...] I think that means a necessary engagement with metaphysical dualism. It is often assumed that hallucinogenic use is very ancient and widespread cross- culturally. Let us call this view the "archaic hallucinogenic use hypothesis." In some passages, you seem to endorse this hypothesis. For example, you note that "[p]sychedelics are associated with pre-modern religious forms and the early history of the current major world religions" (Winkelman, 2017b, p. 2). Now, it would be too long to review the ethnographic and archeological data about hallucinogenic use through history and across cultures, but the example of the Americas should suffice to illustrate why the archaic hallucinogenic use hypothesis is arguably misleading. The New World is generally taken to be the continent where hallucinogenic use has been the most widespread (La Barre, 1964, 1970). So, if the archaic hallucinogenic use hypothesis is accurate, then it should at least be borne out in the Americas; conversely, if it does not accurately depict the Amerindian data, then it is probably wrong everywhere. So, what is the evidence? In South America, we know that ayahuasca use (Brabec de Mori, 2011; Gow, 1992; Shepard, 2014) and jurema use (Samorini, 2018) are relatively recent: probably no more than three or four centuries old. (Note that, following most authors, by "ayahuasca," here, I am referring to the hyperonymic category made of at least Banisteriopsis caapi plus a serotonergic plant (Psychotria viridis or Diplopterys cabrerana) and not to the hyponymic category made of Banisteriopsis caapi alone: the use of the latter is probably more ancient than the use of the former (Miller, Albarracin- Jordan, Moore, \& Capriles, 2019)). We also know that the use of anadenanthera is ancient and was quite widespread in the Caribbean as well as in the Amazon and " Michael Winkelman - The evolutionary neuroanthropology of consciousness 81 ALIUS Bulletin n3 (2019) aliusresearch.org/bulletin Orinoco basins at the time of the arrival of the Europeans (Torres \& Repke, 2006). As regards the oldest use of hallucinogenic snuff, archeological data show that it was confined to very circumscribed places of the central Andes where rather complex cultures have flourished (e.g., Aschero \& Yacobaccio, 1998; Fernandez Distel, 1980). San Pedro also has a long history of use both on the Peruvian coast and in the Andes, but there is no hard archeological data (e.g., chemical analyses) demonstrating that the use of San Pedro is as old as that of anadenanthera (Sharon, 1972). The same goes for brugmansia (toe): it was used in some places in the Andes and in the lowlands, but its use is probably not as old as that of anadenanthera (Gayton, 1928). What about Mesoamerica? The Spanish chroniclers and proto-ethnographers have documented the use of several hallucinogens including various serotonergic mushrooms and peyote (de la Garza, 1990; Heim \& Wasson, 1958, Chapter 1). However, as far as I know, there is no first hand description by any chronicler or proto-ethnographer of any use of these hallucinogenic substances at the East of the Isthmus of Tehuantepec (i.e., in the Mayan lands) (de la Garza, 1990, Chapter 3). All that we have at the East of the Isthmus are somewhat far-fetched interpretations inferring hallucinogenic use from the shape of stones (de Borhegyi, 1961); but no hard evidence whatsoever (Brown, 1984). Moreover, it should be noted that a culture may worship a plant or a mushroom and produce some iconography related to this plant or mushroom without necessarily consuming it or consuming it but at sub-hallucinogenic doses. If we go North, there is evidence of probably several-century-old use of datura in the Southwest (e.g., Stevenson, 1915) but first and foremost in South and Central California (Kroeber, 1976). Nothing more in North America. Admittedly, mescal beans (Troike, 1962) and the "black drink" (Hudson, 1979) were used in the Southern Plains and in the Southeast of North America; however, these substances are psychoactive but not hallucinogenic. Peyote use famously spread from the North of Mexico throughout the Great Plains up to Canada, but this started only at the end of the 19th century (La Barre, 1989). Therefore, it is a very recent phenomenon. Wasson (1979) has argued that some hallucinogens were used in subarctic shamanism. But the only proof he has of it is that today some shamans use Amanita muscaria as a hallucinogen in their rituals. It is quite telling that no proto-ethnographer or explorer has ever reported the use of this mushroom in subarctic Amerindian cultures. It is even more telling that today A. muscaria use among the Ojibwa remains restricted to few shamans and is very controversial (Navet, 2010). All these features betray a recent introduction and use of the fly agaric mushroom. Finally, it must be noted that tobacco has been used in most regions of the Americas-including in the subarctic and circumpolar areas (Winter, 2000b, pp. 9-14)-but only in South America do we find some hallucinogenic use of this plant (Wilbert, 1987). Elsewhere, it was only used at sub-hallucinogenic doses mainly for recreational and psychostimulant purposes (Bollwerk \& Tushingham, 2016; Winter, 2000a). What can be learned from this very sketchy map of hallucinogenic use in the Americas? Well, it looks like hundreds-if not thousands-of Amerindian cultures have not traditionally used any hallucinogens. This is true of almost all cultures of North America (with the exception of some cultures of California and the Southwest) and of all cultures of the South Cone (most notably the Fuegians). It appears that many cultures of the Michael Winkelman - The evolutionary neuroanthropology of consciousness 82 ALIUS Bulletin n3 (2019) aliusresearch.org/bulletin New World did not know hallucinogens, or knew them but were not using them hallucinogenically. This seems to be a problem for the archaic hallucinogenic use hypothesis. By contrast, this is quite consistent with the view advocated by John Cooper (1949) and Johannes Wilbert (1987) according to which hallucinogenic use is closely related to horticultural practices. Do you do endorse a full-fledged version of the "archaic hallucinogenic use hypothesis" or only a weaker version? How do you explain that so many cultures, even in the Americas, have apparently never used hallucinogens? Finally, regarding the etiology of hallucinogenic use, do you side with La Barre (1970) and consider that hallucinogenic use belongs to hunting-based cultures? Or do you alternatively side with Cooper and Wilbert and consider that hallucinogenic use belongs first and foremost to horticulture- based cultures? Well, there is some evidence that visionary substances-I will use this term instead of hallucinogen-increase in importance with horticultural societies. This may reflect increased knowledge, trade networks, etc. or an increased need to adjust to sociocultural change. But just because there is an increase in evidence for the use of visionary plants during periods of social change does not imply that is when the use first started. Where visionary plants are used around the world, they are typically the purview of specialists and often constitute sacred and guarded knowledge. Much of tradition will not be passed on with cultural disintegration-or even development-as may be the case with the famous soma of India. Sometimes the best evidence we have of the ancient entheogenic use of a plant are the local names which involve metaphors alluding to their effects. As to the "hard evidence" for the use of visionary substances in the past, I have pointed to the evidence of enhanced binding of psychedelics with the human serotonergic system (as opposed to other mammals). This is the hard evidence that psychedelics influenced our evolution. Psilocybin-containing mushrooms are the most likely and prevalent psychedelic use in pre-history for a variety of reasons (see my forthcoming "Introduction" to the special issue of the Journal of Psychedelic Studies on "Psychedelics in History and World Religions"). Mushrooms do not leave "hard" evidence unless it is placed into stone or ceramic or metal. And that is indeed what we find around the world, especially throughout Mesoamerica and South America. The presence and placement of mushroom effigies and fungiform representations is too widespread and strategic to ignore their message. Humans everywhere discovered religiosity and spiritual experiences via the impulses that came from psilocybin-containing mushrooms. This is our human spiritual legacy, a deep history of mushroom inspired entheogenic experiences. This impulse has risen Michael Winkelman - The evolutionary neuroanthropology of consciousness 83 ALIUS Bulletin n3 (2019) aliusresearch.org/bulletin and withered many times, but the innate potentials of our brain are always disposed to respond when there is the need and the stimulation. " Humans everywhere discovered religiosity and spiritual experiences via the impulses that came for psilocybin-containing mushrooms. [...]. This impulse has risen and withered many times, but the innate potentials of our brain are always disposed to respond when there is the need and the stimulation. Some scientists study altered states of consciousness in a purely theoretical fashion without exploring those states by themselves. Other scientists care a great deal about having first-hand knowledge of those states before studying them in a third-person fashion. You seem to belong to the second category of scientists. What do you think scientists-and in particular anthropologists-can specifically learn from first-person experience that they cannot learn through third-personal means? The ineffability of mystical and psychedelic experiences is renown. There are things experienced that cannot be completely expressed. It is important to know these kinds of experiences to expand our database and points of reference. Psychedelic experiences are also very personal; the insights obtained are of a personal nature. This kind of knowledge can help us better understand what may and may not be revealed in third person studies. " Science as we know it is largely an ordinary reality construct, a waking mode of consciousness way of ascertaining information about the universe. Psychedelics and other ASCs provide a different epistemic approach [...]. But I think the ultimately most important reason for directly experiencing these substances is what Charles Tart (1972) referred to as "state-specific sciences." Science as we know it is largely an ordinary reality construct, a waking mode of consciousness way of ascertaining information about the universe. Psychedelics and other ASCs provide a different epistemic approach, different ontologies and metaphysics that emerge from the neurophenomenological effects of the substances. We need a science of altered states of consciousness founded in the opportunities " " Michael Winkelman - The evolutionary neuroanthropology of consciousness 84 ALIUS Bulletin n3 (2019) aliusresearch.org/bulletin for knowledge provided by these experiences. Western science is a long way from accomplishing such understandings. Michael Winkelman - The evolutionary neuroanthropology of consciousness 85 ALIUS Bulletin n3 (2019) aliusresearch.org/bulletin References Adams, W., Graf, E., \& Ernst, M. (2004). Experience can change the "light-from-above" prior. Nature Neuroscience, 7(10), 1057-1058. https://doi.org/10.1038/nn1312 Alonso, J. F., Romero, S., Mananas, M. A., \& Riba, J. (2015). Serotonergic psychedelics temporarily modify information transfer in humans. The International Journal of Neuropsychopharmacology, 18(8). https://doi.org/10.1093/ijnp/pyv039 Apicella, C., \& Barrett, H. C. (2016). Cross-cultural evolutionary psychology. Current Opinion in Psychology, 7, 92-97. https://doi.org/10.1016/j.copsyc.2015.08.015 Arthuis, M., Valton, L., Regis, J., Chauvel, P., Wendling, F., Naccache, L., ... Bartolomei, F. (2009). Impaired consciousness during temporal lobe seizures is related to increased long-distance cortical-subcortical synchronization. Brain: A Journal of Neurology, 132, 2091-2101. https://doi.org/10.1093/brain/awp086 Aschero, C. A., \& Yacobaccio, H. D. (1998). 20 Anos Despues: Inca Cueva 7 Reinterpretado. Cuadernos Del Instituto Nacional de Antropologia y Pensamiento Latinoamericano, 18, 7-18. Atasoy, S., Roseman, L., Kaelen, M., Kringelbach, M. L., Deco, G., \& Carhart-Harris, R. L. (2017). Connectome-harmonic decomposition of human brain activity reveals dynamical repertoire re-organization under LSD. Scientific Reports, 7(1), 17661. https://doi.org/10.1038/s41598-017-17546-0 Atran, S. (2002). In Gods We Trust : The Evolutionary Landscape of Religion. New York: Oxford University Press. Barrett, C., \& Kurzban, R. (2006). Modularity in cognition: Framing the debate. Psychological Review, 113(3), 628-647. https://doi.org/10.1037/0033-295X.113.3.628 Barrett, J. L. (2000). Exploring the natural foundations of religion. Trends in Cognitive Sciences, 4(1), 29-34. Barrett, J. L. (2007). Cognitive Science of Religion: What Is It and Why Is It? Religion Compass, 1(6), 768-786. https://doi.org/10.1111/j.1749-8171.2007.00042.x Bayne, T., \& Hohwy, J. (2016). Modes of Consciousness. In W. Sinnott-Armstrong (Ed.), Finding Consciousness: The Neuroscience, Ethics, and Law of Severe Brain Damage. New York: Oxford University Press. Bayne, T., Hohwy, J., \& Owen, A. (2016). Are There Levels of Consciousness? Trends in Cognitive Sciences, 20(6), 405-413. https://doi.org/10.1016/j.tics.2016.03.009 Beggs, J. M., \& Plenz, D. (2003). Neuronal avalanches in neocortical circuits. The Journal of Neuroscience, 23(35), 11167-11177. Michael Winkelman - The evolutionary neuroanthropology of consciousness 86 ALIUS Bulletin n3 (2019) aliusresearch.org/bulletin Billock, V. A., \& Tsou, B. H. (2012). Elementary visual hallucinations and their relationships to neural pattern-forming mechanisms. Psychological Bulletin, 138(4), 744-774. https://doi.org/10.1037/a0027580 Blanke, O., Landis, T., Spinelli, L., \& Seeck, M. (2004). Out-of-body experience and autoscopy of neurological origin. Brain, 127(2), 243-258. https://doi.org/10.1093/brain/awh040 Blumenfeld, H. (2012). Impaired consciousness in epilepsy. The Lancet. Neurology, 11(9), 814-826. https://doi.org/10.1016/S1474-4422(12)70188-6 Bollwerk, E. A., \& Tushingham, S. (Eds.). (2016). Perspectives on the Archaeology of Pipes, Tobacco and other Smoke Plants in the Ancient Americas. Heidelberg: Springer. Boly, M., Perlbarg, V., Marrelec, G., Schabus, M., Laureys, S., Doyon, J., ... Benali, H. (2012). Hierarchical clustering of brain activity during human nonrapid eye movement sleep. Proceedings of the National Academy of Sciences of the United States of America, 109(15), 5856-5861. https://doi.org/10.1073/pnas.1111133109 Bourguignon, E., \& Evascu, T. L. (1977). Altered states of consciousness within a general evolutionary perspective: A holocultural analysis. Behavior Science Research, 12(3), 197-216. https://doi.org/10.1177/106939717701200303 Boyer, P. (2001). Religion Explained: The Evolutionary Origins of Religious Thought. New York: Basic Books. Brabec de Mori, B. (2011). Tracing Hallucinations: Contributing to a Critical Ethnohistory of Ayahuasca Usage in the Peruvian Amazon. In B. Labate \& H. Jungaberle (Eds.), The Internationalization of Ayahuasca (pp. 23-47). Zurich: Lit Verlag. Brown, K. L. (1984). Hallucinogenic Mushrooms, Jade, Obsidian and the Guatemalan Highlands: What Did the Olmecs Really Want? In K. G. Hirth (Ed.), Trade and Exchange in Early Mesoamerica (pp. 215-233). Albuquerque: University of New Mexico Press. Carhart-Harris, R., Erritzoe, D., Williams, T., Stone, J. M., Reed, L. J., Colasanti, A., ... Nutt, D. J. (2012). Neural correlates of the psychedelic state as determined by fMRI studies with psilocybin. Proceedings of the National Academy of Sciences, 109(6), 2138-2143. Carhart-Harris, R. L. (2018). The entropic brain - revisited. Neuropharmacology, 142, 167- 178. https://doi.org/10.1016/j.neuropharm.2018.03.010 Carhart-Harris, R. L., Leech, R., Hellyer, P. J., Shanahan, M., Feilding, A., Tagliazucchi, E., ... Nutt, D. (2014). The entropic brain: A theory of conscious states informed by neuroimaging research with psychedelic drugs. Frontiers in Human Neuroscience, 8. https://doi.org/10.3389/fnhum.2014.00020 Michael Winkelman - The evolutionary neuroanthropology of consciousness 87 ALIUS Bulletin n3 (2019) aliusresearch.org/bulletin Carhart-Harris, R., Muthukumaraswamy, S., Roseman, L., Kaelen, M., Droog, W., Murphy, K., ... Nutt, D. (2016). Neural correlates of the LSD experience revealed by multimodal neuroimaging. Proceedings of the National Academy of Sciences, 113(17), 4853-4858. Chialvo, D., Balenzuela, P., \& Fraiman, D. (2008). The brain: What is critical about it? Conf. Proc. Am. Inst. Phys., 1028, 28-45. Chiao, J., Li, S.-C., Seligman, R., \& Turner, R. (Eds.). (2016). The Oxford handbook of cultural neuroscience. New York: Oxford University Press. Cohen, E. (2007). The Mind Possessed: The Cognition of Spirit Possession in an Afro- Brazilian Religious Tradition. Oxford: Oxford University Press. Cooper, J. M. (1949). Stimulants and narcotics. In J. H. Steward (Ed.), Handbook of South American Indians. Volume 5: The Comparative Ethnology of South American Indians (pp. 525-558). Washington, D.C.: Smithsonian Institution. d'Huy, J. (2013). Le motif du dragon serait paleolithique : mythologie et archeologie. Bulletin Prehistoire Du Sud-Ouest, 21, 195-215. d'Huy, J. (2014). Une methode simple pour reconstruire une mythologie prehistorique (a propos de serpents mythiques sahariens). Cahiers de l'AARS, 17, 95-103. d'Huy, J. (2016). The Evolution of Myths. Scientific American, 315(6), 62-69. https://doi.org/10.1038/scientificamerican1216-62 D'Aquili, E. G., Laughlin, C. D., \& McManus, J. (1979). The Spectrum of Ritual: A Biogenetic Structural Analysis. New York: Columbia University Press. de Borhegyi, S. F. (1961). Miniature Mushroom Stones from Guatemala. American Antiquity, 26(4), 498-504. https://doi.org/10.2307/278737 de la Garza, M. (1990). Sueno y alucinacion en el mundo Nahuatl y Maya. Mexico City: Universidad Nacional Autonoma de Mexico. Descola, P. (2013). Beyond nature and culture (J. Lloyd, Trans.). Chicago: University of Chicago Press. Dokic, J., \& Martin, J.-R. (2012). Disjunctivism, hallucinations, and metacognition. Wiley Interdisciplinary Reviews: Cognitive Science, 3(5), 533-543. Ebert, U., Grossmann, M., Oertel, R., Gramatte, T., \& Kirch, W. (2001). Pharmacokinetic- pharmacodynamic modeling of the electroencephalogram effects of scopolamine in healthy volunteers. Journal of Clinical Pharmacology, 41(1), 51-60. Elman, J., Bates, E., Johnson, M., Karmiloff-Smith, A., Parisi, D., \& Plunkett, K. (1996). Michael Winkelman - The evolutionary neuroanthropology of consciousness 88 ALIUS Bulletin n3 (2019) aliusresearch.org/bulletin Rethinking innateness: A connectionist perspective on development. Cambridge, MA: MIT Press. Fernandez Distel, A. (1980). Hallazgo de pipas en complejos preceramicos del Borde de la Puna Jujena (Republica Argentina) y el empleo de alucinogenos por parte de las mismas culturas. Estudios Arqueologicos, 5, 55-75. Fink, M. (1969). EEG and Human Psychopharmacology. Annual Review of Pharmacology, 9(1), 241-258. https://doi.org/10.1146/annurev.pa.09.040169.001325 Fodor, J. (1983). The Modularity of Mind: An Essay on Faculty Psychology. Cambridge MA: MIT Press. Fortier, M. (In preparation). Manieres neurochimiques de faire des mondes: Neurobiologie, anthropologie et philosophie des experiences hallucinogenes (PhD Dissertation). Institut Jean Nicod, ENS, EHESS, PSL University, Paris. Fortier, M. (2018a). Le sens de realite dans les experiences psychotropes : Etude comparee des hallucinogenes serotoninergiques et anticholinergiques. In S. Baud (Ed.), Histoires et usages des plantes psychotropes (pp. 125-184). Paris: Imago. Fortier, M. (2018b). Sense of reality, metacognition and culture in schizophrenic and drug- induced hallucinations: An interdisciplinary approach. In J. Proust \& M. Fortier (Eds.), Metacognitive Diversity: An Interdisciplinary Approach (pp. 343-378). Oxford/New York: Oxford University Press. Fortier, M. (2019). Le faconnement neuropharmacologique de la culture: Anthropologie comparee des rituels a hallucinogenes serotoninergiques et anticholinergiques. Cahiers d'Anthropologie Sociale, 17, 118-137. Foulkes, W. D. (1962). Dream reports from different stages of sleep. The Journal of Abnormal and Social Psychology, 65(1), 14-25. https://doi.org/10.1037/h0040431 Fox, K. C. R., Dixon, M. L., Nijeboer, S., Girn, M., Floman, J. L., Lifshitz, M., ... Christoff, K. (2016). Functional neuroanatomy of meditation: A review and meta-analysis of 78 functional neuroimaging investigations. Neuroscience and Biobehavioral Reviews, 65, 208-228. https://doi.org/10.1016/j.neubiorev.2016.03.021 Friston, K. J. (2011). Functional and effective connectivity: A review. Brain Connectivity, 1(1), 13-36. https://doi.org/10.1089/brain.2011.0008 Gardner, H. (2000). Intelligence Reframed: Multiple Intelligences for the 21st Century. New York, NY: Basic Books. Gauthier, I., \& Nelson, C. A. (2001). The development of face expertise. Current Opinion in Neurobiology, 11, 219-224. Michael Winkelman - The evolutionary neuroanthropology of consciousness 89 ALIUS Bulletin n3 (2019) aliusresearch.org/bulletin Gayton, A. H. (1928). The Narcotic Plant Datura in Aboriginal American Culture (PhD Dissertation). UC Berkeley, Berkeley, CA. Geisler, W. (2008). Visual perception and the statistical properties of natural scenes. Annual Review of Psychology, 59, 167-192. https://doi.org/10.1146/annurev.psych.58.110405.085632 Geisler, W., \& Kersten, D. (2002). Illusions, perception and Bayes. Nature Neuroscience, 5(6), 508-510. https://doi.org/10.1038/nn0602-508 Gow, P. (1992). River people: Shamanism and history in Western Amazonia. In N. Thomas \& C. Humphrey (Eds.), Shamanism, History, and the State (pp. 90-113). Ann Arbor, MI: University of Michigan Press. Guthrie, S. (1993). Faces in the Clouds: A New Theory of Religion. New York: Oxford University Press. Gyermek, L. (1998). Pharmacology of antimuscarinic agents. Boca Raton FL: CRC Press. Halberstadt, A. L., \& Geyer, M. A. (2011). Multiple receptors contribute to the behavioral effects of indoleamine hallucinogens. Neuropharmacology, 61(3), 364-381. Heim, R., \& Wasson, G. (1958). Les champignons hallucinogenes du Mexique: Etudes ethnologiques, taxinomiques, biologiques, physiologiques et chimiques. Paris: Editions du Museum. Hobson, A. (2001). The dream drugstore: Chemically altered states of consciousness. Cambridge MA: MIT Press. Hudson, C. M. (Ed.). (1979). Black Drink: A Native American Tea. Athens, GA: University of Georgia Press. Ingold, T. (2000). The Perception of the Environment: Essays on livelihood, dwelling and skill. London/New York: Routledge. Itil, T. (1966). Quantitative EEG changes induced by anticholinergic drugs and their behavioral correlates in man. In J. Wortis (Ed.), Recent Advances in Biological Psychiatry: Volume VIII (pp. 151-174). New York: Springer. Itil, T., \& Fink, M. (1966). Anticholinergic drug-induced delirium: Experimental modification, quantitative EEG and behavioral correlations. The Journal of Nervous and Mental Disease, 143(6), 492-507. Jones, D. E. (2016). An Instinct for Dragons. Abingdon: Routledge. Karmiloff-Smith, A. (1992). Beyond modularity: A developmental perspective on cognitive science. Cambridge MA: MIT Press. Michael Winkelman - The evolutionary neuroanthropology of consciousness 90 ALIUS Bulletin n3 (2019) aliusresearch.org/bulletin Karmiloff-Smith, A. (2009). Nativism versus neuroconstructivism: Rethinking the study of developmental disorders. Developmental Psychology, 45(1), 56-63. https://doi.org/10.1037/a0014506 Kent, J. (2010). Psychedelic information theory: Shamanism in the age of reason. Seattle: CreateSpace. Ketchum, J. S. (2006). Chemical warfare: Secrets almost forgotten. A personal story of medical testing of army volunteers with incapacitating chemical agents during the cold war (1955-1975). Santa Rosa, CA: ChemBooks Inc. Ketchum, J. S., Sidell, F., Crowell, E., Aghajanian, G., \& Hayes, A. (1973). Atropine, scopolamine, and Ditran: Comparative pharmacology and antagonists in man. Psychopharmacologia, 28(2), 121-145. Kluver, H. (1966). Mescal and the mechanisms of hallucination. Chicago: University of Chicago Press. Korotayev, A., \& Khaltourina, D. (2011). Myths and genes: A deep historical reconstruction. Moscow: Librokom. Kraehenmann, R. (2017). Dreams and Psychedelics: Neurophenomenological Comparison and Therapeutic Implications. Current Neuropharmacology, 15(7), 1032-1042. https://doi.org/10.2174/1573413713666170619092629 Kraehenmann, R., Pokorny, D., Vollenweider, L., Preller, K. H., Pokorny, T., Seifritz, E., \& Vollenweider, F. X. (2017). Dreamlike effects of LSD on waking imagery in humans depend on serotonin 2A receptor activation. Psychopharmacology, 234(13), 2031- 2046. https://doi.org/10.1007/s00213-017-4610-0 Kroeber, A. (1976). Handbook of the Indians of California. New York: Dover. La Barre, W. (1964). The Narcotic Complex of the New World. Diogenes, 12(48), 125-138. https://doi.org/10.1177/039219216401204808 La Barre, W. (1970). Old and new world narcotics: A statistical question and an ethnological reply. Economic Botany, 24(1), 73-80. https://doi.org/10.1007/BF02860640 La Barre, W. (1989). The peyote cult (Fifth edition, enlarged). Norman/London: University of Oklahoma Press. Laughlin, C. D. (2011). Communing with the Gods: Consciousness, culture and the dreaming brain. Brisbane: Daily Grail. Laughlin, C. D., \& D'Aquili, E. (1974). Biogenetic structuralism. New York: Columbia University Press. Michael Winkelman - The evolutionary neuroanthropology of consciousness 91 ALIUS Bulletin n3 (2019) aliusresearch.org/bulletin Laughlin, C. D., McManus, J., \& D'Aquili, E. (1990). Brain, Symbol and Experience: Toward a Neurophenomenology of Consciousness. Boston MA: Shambhala New Science Library. Le Quellec, J.-L. (2013). Jung et les archetypes: Un mythe contemporain. Auxerre: Editions Sciences Humaines. Lende, D. H., \& Downey, G. (Eds.). (2012). The Encultured Brain: An Introduction to Neuroanthropology. Cambridge MA: MIT Press. Letheby, C. (2016). The Philosophy of Psychedelic Transformation (PhD Dissertation). University of Adelaide, Adelaide. Lewis-Williams, J. D., \& Dowson, T. A. (1988). The Signs of All Times: Entoptic Phenomena in Upper Palaeolithic Art. Current Anthropology, 29(2), 201-245. Luke, D. (2017). Otherworlds: Psychedelics and exceptional human experience. London: Muswell Hill Press. Mandell, A. J. (1980). Toward a psychobiology of transcendence: God in the brain. In J. M. Davidson \& R. J. Davidson (Eds.), The Psychobiology of Consciousness (pp. 379- 464). New York, NY: Plenum Press. Meisel, C., Storch, A., Hallmeyer-Elgner, S., Bullmore, E., \& Gross, T. (2012). Failure of adaptive self-organized criticality during epileptic seizure attacks. PLoS Computational Biology, 8(1), e1002312. https://doi.org/10.1371/journal.pcbi.1002312 Miller, M. J., Albarracin-Jordan, J., Moore, C., \& Capriles, J. M. (2019). Chemical evidence for the use of multiple psychotropic plants in a 1,000-year-old ritual bundle from South America. Proceedings of the National Academy of Sciences, 1-6. https://doi.org/10.1073/pnas.1902174116 Milliere, R., Carhart-Harris, R. L., Roseman, L., Trautwein, F.-M., \& Berkovich-Ohana, A. (2018). Psychedelics, Meditation, and Self-Consciousness. Frontiers in Psychology, 9, 1-29. https://doi.org/10.3389/fpsyg.2018.01475 Murdock, G. P., Wilson, S. F., \& Frederick, V. (1980). World Distribution of Theories of Illness. Transcultural Psychiatric Research Review, 17(1-2), 37-64. https://doi.org/10.1177/136346158001700102 Muthukumaraswamy, S., Carhart-Harris, R., Moran, R., Brookes, M., Williams, T., Erritzoe, D., ... Nutt, D. (2013). Broadband cortical desynchronization underlies the human psychedelic state. Journal of Neuroscience, 33(38), 15171-15183. https://doi.org/10.1523/JNEUROSCI.2063-13.2013 Navet, E. (2010). Les voies de l'amanite tue-mouches (Amanita muscaria): entre "enfer" et "paradis." In S. Baud \& C. Ghasarian (Eds.), Des plantes psychotropes: Initiations, Michael Winkelman - The evolutionary neuroanthropology of consciousness 92 ALIUS Bulletin n3 (2019) aliusresearch.org/bulletin therapies et quetes de soi (pp. 61-78). Paris: IMAGO. Nir, Y., \& Tononi, G. (2010). Dreaming and the brain: From phenomenology to neurophysiology. Trends in Cognitive Sciences, 14(2), 88. https://doi.org/10.1016/j.tics.2009.12.001 Osipova, D., Ahveninen, J., Kaakkola, S., Jaaskelainen, I. P., Huttunen, J., \& Pekkonen, E. (2003). Effects of scopolamine on MEG spectral power and coherence in elderly subjects. Clinical Neurophysiology: Official Journal of the International Federation of Clinical Neurophysiology, 114(10), 1902-1907. Palhano-Fontes, F., Andrade, K. C., Tofoli, L. F., Santos, A. C., Crippa, J. A. S., Hallak, J. E. C., ... Araujo, D. B. de. (2015). The psychedelic state induced by ayahuasca modulates the activity and connectivity of the default mode network. PLOS ONE, 10(2), e0118143. Panel on Anticholinesterase Chemicals, Panel on Anticholinergic Chemicals, Committee on Toxicology, Board on Toxicology and Environmental Health Hazards, \& Assembly of Life Sciences. (1982). Possible long-term health effects of short-term exposure to chemical agents. Volume 1: Anticholinesterases and anticholinergics (Electronic version). Washington, DC: National Academy Press. Previc, F. H. (2006). The role of the extrapersonal brain systems in religious activity. Consciousness and Cognition, 15(3), 500-539. https://doi.org/10.1016/j.concog.2005.09.009 Previc, F. H. (2009). The Dopaminergic Mind in Human Evolution and History. Cambridge: Cambridge University Press. Priesemann, V., Valderrama, M., Wibral, M., \& Le Van Quyen, M. (2013). Neuronal Avalanches Differ from Wakefulness to Deep Sleep - Evidence from Intracranial Depth Recordings in Humans. PLOS Computational Biology, 9(3), e1002985. https://doi.org/10.1371/journal.pcbi.1002985 Pyysiainen, I. (2013). Cognitive Science of Religion: State-of-the-Art. Journal for the Cognitive Science of Religion, 1(1), 5-28. Ray, T. S. (2010). Psychedelics and the Human Receptorome. PLoS ONE, 5(2). https://doi.org/10.1371/journal.pone.0009019 Ray, T. S. (2013). Mental Organs and the Origins of Mind. In L. Swan (Ed.), Origins of Mind (pp. 301-326). Heidelberg/New York: Springer. Ray, T. S. (2016). Constructing the ecstasy of MDMA from its component mental organs: Proposing the primer/probe method. Medical Hypotheses, 87, 48-60. https://doi.org/10.1016/j.mehy.2015.12.018 Michael Winkelman - The evolutionary neuroanthropology of consciousness 93 ALIUS Bulletin n3 (2019) aliusresearch.org/bulletin Roepstorff, A., Niewohner, J., \& Beck, S. (2010). Enculturing brains through patterned practices. Neural Networks: The Official Journal of the International Neural Network Society, 23(8-9), 1051-1059. Roseman, L., Leech, R., Feilding, A., Nutt, D. J., \& Carhart-Harris, R. L. (2014). The effects of psilocybin and MDMA on between-network resting state functional connectivity in healthy volunteers. Frontiers in Human Neuroscience, 8. https://doi.org/10.3389/fnhum.2014.00204 Samorini, G. (2018). El culto a la Jurema en los documentos historicos. Revista Cultura y Droga, 23(25), 33-47. https://doi.org/10.17151/culdr.2018.23.25. Sanz, C., Zamberlan, F., Erowid, E., Erowid, F., \& Tagliazucchi, E. (2018). The Experience Elicited by Hallucinogens Presents the Highest Similarity to Dreaming within a Large Database of Psychoactive Substance Reports. Frontiers in Neuroscience, 12. https://doi.org/10.3389/fnins.2018.00007 Schenberg, E. E., Alexandre, J. F. M., Filev, R., Cravo, A. M., Sato, J. R., Muthukumaraswamy, S., ... Silveira, D. X. da. (2015). Acute Biphasic Effects of Ayahuasca. PLOS ONE, 10(9). https://doi.org/10.1371/journal.pone.0137202 Sharon, D. (1972). The San Pedro Cactus in Peruvian Folk Healing. In P. T. Furst (Ed.), Flesh of the Gods: The Ritual Use of Hallucinogens (pp. 114-135). New York/Washington: Praeger. Shepard, G. (2014). Will the real shaman please stand up? The recent adoption of ayahuasca among indigenous groups of the Peruvian Amazon. In B. Labate \& C. Cavnar (Eds.), Ayahuasa shamanism in the Amazon and beyond (pp. 16-39). Oxford/New York: Oxford University Press. Siclari, F., Baird, B., Perogamvros, L., Bernardi, G., LaRocque, J. J., Riedner, B., ... Tononi, G. (2017). The neural correlates of dreaming. Nature Neuroscience, 20(6), 872-878. https://doi.org/10.1038/nn.4545 Sosis, R. (2009). The Adaptationist-Byproduct Debate on the Evolution of Religion: Five Misunderstandings of the Adaptationist Program. Journal of Cognition and Culture, 9(3-4), 315-332. https://doi.org/10.1163/156770909X12518536414411 Sperber, D. (1996a). Explaining Culture: A Naturalistic Approach. Oxford/Malden MA: Blackwell. Sperber, D. (1996b). Why are perfect animals, hybrids, and monsters food for symbolic thought? Method \& Theory in the Study of Religion, 8(2), 143-169. Sperber, D., \& Hirschfeld, L. A. (2004). The cognitive foundations of cultural stability and diversity. Trends in Cognitive Sciences, 8(1), 40-46. https://doi.org/10.1016/j.tics.2003.11.002 Michael Winkelman - The evolutionary neuroanthropology of consciousness 94 ALIUS Bulletin n3 (2019) aliusresearch.org/bulletin Stevenson, M. C. (1915). Ethnobotany of the Zuni Indians. In Thirtieth Annual Report of the Bureau of American Ethnology, 1908-1909 (pp. 35-102). Washington, D.C.: Smithsonian Institution. Stickgold, R., Malia, A., Fosse, R., Propper, R., \& Hobson, J. A. (2001). Brain-mind states: I. Longitudinal field study of sleep/wake factors influencing mentation report length. Sleep, 24(2), 171-179. https://doi.org/10.1093/sleep/24.2.171 Strassman, R., Wojtowicz, S., Luna, L. E., \& Frecska, E. (2008). Inner paths to outer space: Journeys to alien worlds through psychedelics and other spiritual technologies. Rochester VT: Park Street Press. Sun, J., \& Perona, P. (1998). Where is the sun? Nature Neuroscience, 1(3), 183-184. https://doi.org/10.1038/630 Tagliazucchi, E. (2017). The signatures of conscious access and its phenomenology are consistent with large-scale brain communication at criticality. Consciousness and Cognition, 55, 136-147. https://doi.org/10.1016/j.concog.2017.08.008 Tagliazucchi, E., \& Chialvo, D. R. (2013). Brain complexity born out of criticality. ArXiv, 4-13. https://doi.org/10.1063/1.4776495 Tagliazucchi, E., von Wegner, F., Morzelewski, A., Brodbeck, V., Borisov, S., Jahnke, K., \& Laufs, H. (2013). Large-scale brain functional modularity is reflected in slow electroencephalographic rhythms across the human non-rapid eye movement sleep cycle. NeuroImage, 70, 327-339. https://doi.org/10.1016/j.neuroimage.2012.12.073 Tart, C. T. (1972). States of Consciousness and State-Specific Sciences. Science, 176(4040), 1203-1210. https://doi.org/10.1126/science.176.4040.1203 Testart, A. (2012). Avant l'histoire: L'evolution des societes, de Lascaux a Carnac. Paris: Gallimard. Timme, N. M., Marshall, N. J., Bennett, N., Ripp, M., Lautzenhiser, E., \& Beggs, J. M. (2016). Criticality Maximizes Complexity in Neural Tissue. Frontiers in Physiology, 7. https://doi.org/10.3389/fphys.2016.00425 Tononi, G., Edelman, G. M., \& Sporns, O. (1998). Complexity and coherency: Integrating information in the brain. Trends in Cognitive Sciences, 2(12), 474-484. Tononi, G., Sporns, O., \& Edelman, G. M. (1996). A complexity measure for selective matching of signals by the brain. Proceedings of the National Academy of Sciences, 93(8), 3422-3427. https://doi.org/10.1073/pnas.93.8.3422 Tooby, J., \& Cosmides, L. (1992). The psychological foundations of culture. In J. Barkow, L. Cosmides, \& J. Tooby (Eds.), The adapted mind: Evolutionary psychology and the generation of culture (pp. 19-136). Oxford: Oxford University Press. Michael Winkelman - The evolutionary neuroanthropology of consciousness 95 ALIUS Bulletin n3 (2019) aliusresearch.org/bulletin Torres, C. M., \& Repke, D. B. (2006). Anadenanthera: Visionary plant of ancient South America. New York: Haworth Press. Troike, R. C. (1962). The Origins of Plains Mescalism. American Anthropologist, 64(5), 946-963. https://doi.org/10.1525/aa.1962.64.5.02a00040 Tylor, E. (1871). Primitive culture: Researches into the development of mythology, philosophy, religion, language, art and custom. Londres: John Murray. Vollenweider, Franz X. (1998). Recent advances and concepts in the search for biological correlates of hallucinogen-induced altered states of consciousness. The Heffter Review of Psychedelic Research, 1, 21-32. Vollenweider, Franz X. (2001). Brain mechanisms of hallucinogens and entactogens. Dialogues in Clinical Neuroscience, 3(4), 265-279. Vollenweider, Franz X., \& Geyer, M. A. (2001). A systems model of altered consciousness: Integrating natural and drug-induced psychoses. Brain Research Bulletin, 56(5), 495-507. Wasson, R. G. (1979). Traditional Use in North America of Amanita muscaria for Divinatory Purposes. Journal of Psychedelic Drugs, 11(1-2), 25-28. https://doi.org/10.1080/02791072.1979.10472088 Wengrow, D. (2013). The Origins of Monsters: Image and Cognition in the First Age of Mechanical Reproduction. Princeton, NJ: Princeton University Press. Wilbert, J. (1987). Tobacco and Shamanism in South America. New Haven: Yale University Press. Winkelman, M. (1982). Magic: A Theoretical Reassessment. Current Anthropology, 23(1), 37-66. Winkelman, M. (1986). Magico-religious practitioner types and socioeconomic conditions. Cross-Cultural Research, 20(1-4), 17-46. Winkelman, M. (1990). Shamans and Other "Magico-Religious" Healers: A Cross-Cultural Study of Their Origins, Nature, and Social Transformations. Ethos, 18(3), 308-352. Winkelman, M. (1992). Shamans, priests, and witches: A cross-cultural study of magico- religious practitioners. Tempe: Arizona State University. Winkelman, M. (1996). Psychointegrator plants: Their roles in human culture and health. In M. Winkelman \& W. Andritzky (Eds.), Sacred plants, consciousness, and healing: Cross-cultural and interdisciplinary perspectives (pp. 9-53). Berlin: Verlag fur Wissenschaft und Bildung. Michael Winkelman - The evolutionary neuroanthropology of consciousness 96 ALIUS Bulletin n3 (2019) aliusresearch.org/bulletin Winkelman, M. (2001). Psychointegrators: Multidisciplinary Perspectives on the Therapeutic Effects of Hallucinogens. Complementary Health Practice Review, 6(3), 219-237. https://doi.org/10.1177/153321010100600304 Winkelman, M. (2002a). Shamanism and Cognitive Evolution. Cambridge Archaeological Journal, 12(1), 71-101. https://doi.org/10.1017/S0959774302000045 Winkelman, M. (2002b). Shamanism as Neurotheology and Evolutionary Psychology. American Behavioral Scientist, 45(12), 1873-1885. Winkelman, M. (2007). Therapeutic bases of psychedelic medicines: Psychointegrative effects. In M. Winkelman \& T. B. Roberts (Eds.), Psychedelic medicine: New evidence for hallucinogenic substances as treatments (pp. 1-19). Westport, CT: Praeger. Winkelman, M. (2009). Culture and Health: Applying Medical Anthropology. San Francisco, CA: Jossey-Bass. Winkelman, M. (2010a). Shamanism: A Biopsychosocial Paradigm of Consciousness and Healing (2d edition). Santa Barbara: Praeger. Winkelman, M. (2010b). Shamanism and the Origins of Spirituality and Ritual Healing. Journal for the Study of Religion, Nature and Culture, 3(4), 458-489. https://doi.org/10.1558/jsrnc.v3i4.458 Winkelman, M. (2010c). The Shamanic Paradigm: Evidence from Ethnology, Neuropsychology and Ethology. Time and Mind, 3(2), 159-181. https://doi.org/10.2752/175169610X12632240392758 Winkelman, M. (2011). A paradigm for understanding altered consciousness: The integrative mode of consciousness. In E. Cardena \& M. Winkelman (Eds.), Altering consciousness: Multidisciplinary perspectives. Volume 1: History, culture, and the humanities (pp. 23-41). Santa Barbara: Praeger. Winkelman, M. (2013). Shamanic Cosmology as an Evolutionary Neurocognitive Epistemology. International Journal of Transpersonal Studies, 32(1). https://doi.org/10.24972/ijts.2013.32.1.79 Winkelman, M. (2017a). Shamanism and the brain. In N. K. Clements (Ed.), Religion: Mental religion (pp. 355-372). New York, NY: MacMillan. Winkelman, M. (2017b). The Mechanisms of Psychedelic Visionary Experiences: Hypotheses from Evolutionary Psychology. Frontiers in Neuroscience, 11. https://doi.org/10.3389/fnins.2017.00539 Winkelman, M. (2018a). An ontology of psychedelic entity experiences in evolutionary psychology and neurophenomenology. Journal of Psychedelic Studies, 2(1), 5-23. Michael Winkelman - The evolutionary neuroanthropology of consciousness 97 ALIUS Bulletin n3 (2019) aliusresearch.org/bulletin https://doi.org/10.1556/2054.2018.002 Winkelman, M. (2018b). Shamanism and Possession. In The International Encyclopedia of Anthropology. https://doi.org/10.1002/9781118924396.wbiea1651 Winkelman, M., \& Baker, J. R. (2010). Supernatural as natural: A biocultural approach to religion. Upper Saddle River, NJ: Prentice Hall. Winter, J. C. (Ed.). (2000a). Tobacco Use by Native North Americans: Sacred Smoke and Silent Killer. Norman, OK: University of Oklahoma Press. Winter, J. C. (2000b). Traditional Uses of Tobacco by Native Americans. In J. C. Winter (Ed.), Tobacco Use by Native Americans: Sacred Smoke and Silent Killer (pp. 9- 58). Norman, OK: University of Oklahoma Press. Yang, Y., Cheng, X., Liu, W., Chou, G., Wang, Z., \& Wang, C. (2015). Potent AChE and BChE inhibitors isolated from seeds of Peganum harmala Linn by a bioassay-guided fractionation. Journal of Ethnopharmacology, 168, 279-286. https://doi.org/10.1016/j.jep.2015.03.070 Zhao, T., Ding, K., Zhang, L., Cheng, X., Wang, C., \& Wang, Z. (2013). Acetylcholinesterase and Butyrylcholinesterase Inhibitory Activities of -Carboline and Quinoline Alkaloids Derivatives from the Plants of Genus Peganum. Journal of Chemistry. https://doi.org/dx.doi.org/10.1155/2013/717232 Zheng, X., Zhang, Z., Chou, G., Wu, T., Cheng, X., Wang, C., \& Wang, Z. (2009). Acetylcholinesterase inhibitive activity-guided isolation of two new alkaloids from seeds of Peganum nigellastrum Bunge by an in vitro TLC- bioautographic assay. Archives of Pharmacal Research, 32(9), 1245-1251. https://doi.org/10.1007/s12272- 009-1910-x\par
  \endgroup
}%
\end{document}
